% =============================================================================
%  Notas de Análisis I — MA0505 — III Parcial
%  Compilar con: pdflatex Analisis_I_ExamenIII.tex (dos veces para que el TOC
%  recoja las entradas de cada definición, teorema, lema, etc.)
% =============================================================================
\documentclass[11pt,twoside]{article}

\usepackage[margin=0.9in]{geometry}
\usepackage[utf8]{inputenc}
\usepackage[T1]{fontenc}
\usepackage{XCharter}
\usepackage[libertine,charter]{newtxmath}
\usepackage[spanish,es-noshorthands]{babel}
% Renombramos elementos al español manualmente (no requiere texlive-lang-spanish)
\addto\captionsenglish{%
  \renewcommand{\contentsname}{\'Indice}%
  \renewcommand{\refname}{Referencias}%
}
\usepackage{amsmath,amssymb,amsthm,mathtools}
\usepackage{enumitem}
\usepackage[protrusion=true,expansion=false]{microtype}
\usepackage{array}
\usepackage{longtable}
\usepackage{titlesec}
\usepackage{xcolor}
\usepackage[most]{tcolorbox}
\usepackage{tikz}
\usetikzlibrary{positioning, arrows.meta, shapes.geometric, calc, fit}
\usepackage{eso-pic}
\usepackage{etoolbox}
\usepackage{xstring}
\usepackage{fancyhdr}
\usepackage{lastpage}
\usepackage{hyperref}

\hypersetup{
  colorlinks=true,
  linkcolor=blue!50!black,
  urlcolor=blue!50!black,
  pdftitle={Notas de Análisis I --- MA0505},
  pdfauthor={Sebastián Fernández Rivera}
}

% Mostrar entradas hasta nivel "subsubsection" en el índice
% (cada teorema/definición/etc. aparece en el TOC con nombre y página).
\setcounter{tocdepth}{3}
\setcounter{secnumdepth}{2}

\renewcommand{\proofname}{\textbf{Prueba}}
\allowdisplaybreaks

% -----------------------------------------------------------------------------
%  Contador y entornos para resultados (definiciones, teoremas, etc.)
%  - Cada resultado se numera dentro de su subsection: 1.1.1, 1.1.2, ...
%  - Cada resultado se agrega automáticamente al índice con su nombre y página.
% -----------------------------------------------------------------------------
\newcounter{result}[subsection]
\renewcommand{\theresult}{\thesubsection.\arabic{result}}

% Cuerpo en cursiva (para teoremas, proposiciones, lemas, corolarios)
\newenvironment{namedres}[2]{%
  \refstepcounter{result}%
  \phantomsection%
  \addcontentsline{toc}{subsubsection}{\protect\numberline{\theresult}#1 (\textit{#2})}%
  \par\medskip\noindent\textbf{#1 \theresult\ (\textit{#2}).}%
  \par\nopagebreak\noindent\itshape\ignorespaces%
}{\par\medskip\upshape}

% -----------------------------------------------------------------------------
%  Caja coloreada para los resúmenes de los anexos.
%  Uso: \begin{equivbox}{ColorXcolor}{Título}  ...  \end{equivbox}
% -----------------------------------------------------------------------------
\newtcolorbox{equivbox}[2]{%
  enhanced,
  breakable,
  colback=#1!6!white,
  colframe=#1!72!black,
  coltitle=white,
  fonttitle=\bfseries,
  fontupper=\small,
  title={#2},% braces protect commas in the title from tcolorbox key parser
  arc=4pt,
  boxrule=0.8pt,
  left=10pt, right=10pt, top=6pt, bottom=6pt,
  before skip=4pt, after skip=8pt,
}

% Encabezado compacto para los sub-grupos de un Anexo:
% sustituye a \subsection*{...}\addcontentsline{...} cuando se desea poco
% espacio entre el título del grupo y la caja siguiente.
\newcommand{\anexogroup}[1]{%
  \par\addvspace{0.7\baselineskip}%
  \noindent{\bfseries\large #1}\par\nobreak
  \addcontentsline{toc}{subsection}{#1}%
  \vspace{4pt}\nopagebreak
}

% Cuerpo en redonda (para definiciones, ejemplos, ejercicios, notas, notación)
\newenvironment{namedstd}[2]{%
  \refstepcounter{result}%
  \phantomsection%
  \addcontentsline{toc}{subsubsection}{\protect\numberline{\theresult}#1 (\textit{#2})}%
  \par\medskip\noindent\textbf{#1 \theresult\ (\textit{#2}).}%
  \par\nopagebreak\noindent\ignorespaces%
}{\par\medskip}

% Subrayado HTML <u>...</u> -> \uline
\newcommand{\hlu}[1]{\underline{#1}}

% Operadores y abreviaturas
\DeclareMathOperator{\Ker}{Ker}
\DeclareMathOperator{\Imm}{Im}
\DeclareMathOperator{\Aut}{Aut}
\DeclareMathOperator{\sgn}{sgn}
\DeclareMathOperator{\mcd}{mcd}
\DeclareMathOperator{\MCD}{MCD}
\newcommand{\Z}{\mathbb{Z}}
\newcommand{\R}{\mathbb{R}}
\newcommand{\Q}{\mathbb{Q}}
\newcommand{\N}{\mathbb{N}}
\newcommand{\F}{\mathbb{F}}
% Función indicadora (los beamers usan \ind)
\newcommand{\ind}{\mathbf{1}}

% =============================================================================
%  Cabecera y pie de página (fancyhdr)
% =============================================================================
\pagestyle{fancy}
\fancyhf{}
\renewcommand{\sectionmark}[1]{\markboth{\thesection.\ #1}{}}
\fancyhead[L]{\small\nouppercase{\leftmark}}
\fancyhead[R]{\small\textit{MA0505 --- An\'alisis I, III Parcial}}
\fancyfoot[C]{\small Página \thepage{} de \pageref{LastPage}}
\renewcommand{\headrulewidth}{0.4pt}
\renewcommand{\footrulewidth}{0pt}
\setlength{\headheight}{14pt}

% =============================================================================
%  Tabs de color en el borde exterior (índice rápido por sección)
% =============================================================================
\definecolor{sec1col}{HTML}{1F77B4}   % Egorov y Lusin
\definecolor{sec2col}{HTML}{2CA02C}   % Convergencia en medida
\definecolor{sec3col}{HTML}{17BECF}   % Integral de funciones no negativas
\definecolor{sec4col}{HTML}{9467BD}   % Propiedades de la integral
\definecolor{sec5col}{HTML}{FF7F0E}   % Integral de funciones medibles
\definecolor{sec6col}{HTML}{D62728}   % Riemann y Lebesgue
\definecolor{sec7col}{HTML}{8C564B}   % (reserva)
\definecolor{sec8col}{HTML}{E377C2}   % (reserva)
\definecolor{sec9col}{HTML}{BCBD22}   % (reserva)
\definecolor{sec10col}{HTML}{555555}  % Anexo

\newcounter{thumbsec}
\pretocmd{\section}{\stepcounter{thumbsec}}{}{}

% Bandera para evitar que los thumb-tabs aparezcan en la portada y el índice.
% Se activa explícitamente en el cuerpo, justo después del \tableofcontents.
\newif\ifdrawthumb
\drawthumbfalse

\newcommand{\drawthumbtab}{%
  \ifdrawthumb%
  \ifnum\value{thumbsec}>0%
    \ifnum\value{thumbsec}<7%
      \begin{tikzpicture}[remember picture, overlay]%
        \pgfmathsetmacro{\toppadcm}{4.0}%
        \pgfmathsetmacro{\tabHcm}{1.6}%
        \pgfmathsetmacro{\ypos}{-\toppadcm - (\value{thumbsec} - 1) * \tabHcm}%
        \ifodd\value{page}%
          % Páginas impares (recto) --- tab en el borde derecho.
          \fill[sec\arabic{thumbsec}col]
            ([xshift=0pt, yshift=\ypos cm]current page.north east)
            rectangle ++(-0.7cm, -1.2cm);
          \node[white, font=\bfseries\large, anchor=center]
            at ([xshift=-0.35cm, yshift=\ypos cm - 0.6cm]current page.north east)
            {\arabic{thumbsec}};
        \else%
          % Páginas pares (verso) --- tab en el borde izquierdo.
          \fill[sec\arabic{thumbsec}col]
            ([xshift=0pt, yshift=\ypos cm]current page.north west)
            rectangle ++(0.7cm, -1.2cm);
          \node[white, font=\bfseries\large, anchor=center]
            at ([xshift=0.35cm, yshift=\ypos cm - 0.6cm]current page.north west)
            {\arabic{thumbsec}};
        \fi%
      \end{tikzpicture}%
    \fi%
  \fi%
  \fi%
}
\AddToShipoutPictureFG{\drawthumbtab}

\begin{document}

% =============================================================================
% PORTADA
% =============================================================================
\begin{titlepage}
\thispagestyle{empty}
\centering

\vspace*{1.5cm}

{\scshape\Large Universidad de Costa Rica\par}
\vspace{0.3cm}
{\scshape\large Escuela de Matem\'atica\par}

\vspace{1.2cm}

{\color{blue!50!black}\rule{\textwidth}{1.2pt}}
\vspace{0.6cm}

{\Huge\bfseries Notas de An\'alisis I\par}

\vspace{0.5cm}

{\Large\itshape Resumen de definiciones, teoremas y resultados\par}

\vspace{0.4cm}
{\color{blue!50!black}\rule{\textwidth}{1.2pt}}

\vspace{1.4cm}

{\Large\bfseries MA0505 --- An\'alisis I\par}

\vspace{2.2cm}

\begin{tabular}{rl}
\textbf{Autor:}              & Sebasti\'{a}n Fern\'{a}ndez Rivera \\[4pt]
\textbf{Curso:}              & MA0505 --- An\'alisis I \\[4pt]
\textbf{Documento:}          & Notas para el III Parcial \\[4pt]
\textbf{Resultados numerados:} & 63 \\[4pt]
\textbf{Fecha:}              & \today
\end{tabular}

\vfill

\begin{minipage}{0.85\textwidth}
\centering
\itshape
Este documento re\'{u}ne las definiciones, teoremas, proposiciones, lemas, corolarios, ejemplos, notas, ejercicios y notaciones presentadas en el curso sobre los teoremas de Egorov y Lusin, la convergencia en medida, la integral de Lebesgue y su relaci\'{o}n con la integral de Riemann. Cada resultado est\'{a} numerado de la forma \texttt{seccion.subseccion.resultado} y aparece en el \'{\i}ndice con su p\'{a}gina correspondiente para una consulta r\'{a}pida durante el examen.
\end{minipage}

\vspace{1cm}

{\color{blue!50!black}\rule{0.5\textwidth}{0.6pt}}

\end{titlepage}

% =============================================================================
% \'INDICE
% =============================================================================
\thispagestyle{empty}
{\hypersetup{linkcolor=black}\tableofcontents}
\newpage
% \tableofcontents llama internamente a \section*{...}, lo cual activa el hook
% \pretocmd{\section}{...} y hace avanzar \value{thumbsec} a 1 antes de iniciar
% el cuerpo. Reseteamos para que la sección 1 reciba el thumb-tab "1".
\setcounter{thumbsec}{0}
\drawthumbtrue   % Activa los thumb-tabs a partir de aquí (cuerpo del documento).

% =============================================================================
\section{Los teoremas de Egorov y Lusin}
% =============================================================================

\subsection{El teorema de Egorov}

A lo largo de esta sección, $m$ denota la medida de Lebesgue en $\R^{d}$, $m_{e}$ la medida exterior, y ``c.p.d.'' abrevia \emph{casi por doquier}.

\begin{namedstd}{Nota}{La convergencia puntual no implica convergencia uniforme}
Recordemos que existen sucesiones $f_{k} : E \to \R$ continuas tales que
\[
\lim_{k \to \infty} f_{k} = f
\]
puntualmente en $E$, pero tales que $f_{k}$ no converge uniformemente a $f$ en compactos. El teorema de Egorov muestra que, para funciones medibles sobre un conjunto de medida finita, la convergencia puntual c.p.d. sí es uniforme fuera de un conjunto de medida arbitrariamente pequeña.
\end{namedstd}

\begin{namedstd}{Ejemplo}{La hipótesis de medida finita es esencial en Egorov}
Sean $E = \R^{d}$ y
\[
f_{k}(x) = \ind_{B(0,k)}(x).
\]
Entonces $\lim_{k \to \infty} f_{k} = 1$ puntualmente en $\R^{d}$. Sea $F$ cerrado y no acotado; entonces para todo $k \in \N$ existe $x \in F$ tal que
\[
|1 - f_{k}(x)| = 1,
\]
pues $F$ posee puntos fuera de $B(0,k)$. Es decir, la convergencia no puede ser uniforme en $F$. Note además que, si $F$ es cerrado y $m(\R^{d} \setminus F) < \infty$, entonces $F$ no es acotado. Por lo tanto, la conclusión del teorema de Egorov falla para esta sucesión: ningún cerrado que deje afuera medida finita admite convergencia uniforme.
\end{namedstd}

\begin{namedres}{Lema}{Lema técnico previo al teorema de Egorov}
Sean $E \subseteq \R^{d}$ medible con $m(E) < \infty$ y $\{f_{k}\}_{k=1}^{\infty}$ funciones medibles en $E$ tales que
\[
\lim_{k \to \infty} f_{k} = f \quad \text{c.p.d. en } E, \qquad |f(x)| \neq \infty \ \text{c.p.d. en } E.
\]
Entonces, dados $\varepsilon > 0$ y $\eta > 0$, existen $F \subseteq E$ cerrado y $k_{0} = k_{0}(\varepsilon, \eta) \geq 0$ tales que
\[
m(E \setminus F) < \eta
\]
y
\[
|f_{k}(x) - f(x)| < \varepsilon \quad \text{para todo } x \in F \text{ y } k \geq k_{0}.
\]
\end{namedres}
\begin{proof}
Defina
\[
\tilde{E} = \{ x \in E :\ \lim_{k \to \infty} f_{k}(x) = f(x),\ |f(x)| < \infty \}.
\]
Por hipótesis, $m(E \setminus \tilde{E}) = 0$. Para $\varepsilon > 0$ y $\eta > 0$ fijos, definamos
\[
E_{m} = \{ x \in \tilde{E} :\ k \geq m \implies |f(x) - f_{k}(x)| < \varepsilon \}.
\]
Entonces cada $E_{m}$ es medible, pues
\[
E_{m} = \tilde{E} \cap \bigcap_{k \geq m} \{ x \in E :\ |f(x) - f_{k}(x)| < \varepsilon \}
\]
es una intersección numerable de conjuntos medibles, y además $E_{m} \subseteq E_{m+1}$, ya que la condición se exige sobre menos índices $k$ al crecer $m$. Se cumple que $\bigcup_{m=1}^{\infty} E_{m} = \tilde{E}$ (ejercicio; véase el ejercicio siguiente).

Por la continuidad desde abajo de la medida, tenemos que
\[
m(\tilde{E}) = \lim_{m \to \infty} m(E_{m}).
\]
Como $m(\tilde{E}) \leq m(E) < \infty$, podemos restar y obtenemos
\[
\lim_{m \to \infty} m(\tilde{E} \setminus E_{m}) = \lim_{m \to \infty} \big( m(\tilde{E}) - m(E_{m}) \big) = 0.
\]
Además $E \setminus E_{m} = (E \setminus \tilde{E}) \cup (\tilde{E} \setminus E_{m})$ y $m(E \setminus \tilde{E}) = 0$, entonces $m(E \setminus E_{m}) = m(\tilde{E} \setminus E_{m}) \to 0$. Sea entonces $k_{0}$ tal que
\[
m(E \setminus E_{k_{0}}) < \frac{\eta}{2}.
\]
Note que si $x \in E_{k_{0}}$, entonces, por definición de $E_{k_{0}}$,
\[
|f(x) - f_{k}(x)| < \varepsilon \quad \text{para } k \geq k_{0}.
\]
Finalmente, como $E_{k_{0}}$ es medible, por regularidad interna existe $F$ cerrado con $F \subseteq E_{k_{0}}$ que satisface
\[
m(E_{k_{0}} \setminus F) < \frac{\eta}{2}.
\]
Entonces
\[
m(E \setminus F) \leq m(E \setminus E_{k_{0}}) + m(E_{k_{0}} \setminus F) < \eta,
\]
y si $x \in F \subseteq E_{k_{0}}$ y $k \geq k_{0}$, tenemos que $|f_{k}(x) - f(x)| < \varepsilon$.
\end{proof}

\begin{namedstd}{Ejercicio}{La unión de los $E_{m}$ recupera $\tilde{E}$}
Con la notación de la prueba anterior, muestre que $\bigcup_{m=1}^{\infty} E_{m} = \tilde{E}$.
\end{namedstd}

\begin{namedres}{Teorema}{Egorov: la convergencia c.p.d. es casi uniforme en medida finita}
Sean $E \subseteq \R^{d}$ medible, con $m(E) < \infty$, y $\{f_{k}\}_{k=1}^{\infty}$ funciones medibles en $E$. Asuma que
\[
\lim_{k \to \infty} f_{k} = f
\]
c.p.d. en $E$ y que $|f(x)| \neq \infty$ c.p.d. en $E$. Entonces, dado $\varepsilon > 0$, existe $F_{\varepsilon} \subseteq E$ cerrado tal que:
\begin{enumerate}[label=(\roman*)]
\item $m(E \setminus F_{\varepsilon}) < \varepsilon$.
\item $f_{k}$ converge a $f$ uniformemente en $F_{\varepsilon}$.
\end{enumerate}
\end{namedres}
\begin{proof}
Sea $\varepsilon > 0$. Aplicando el lema anterior con $\varepsilon_{m} = \tfrac{1}{m}$ y $\eta_{m} = \tfrac{\varepsilon}{2^{m}}$ para cada $m \geq 1$, existen $F_{m} \subseteq E$ cerrados y enteros $\kappa_{m}^{\varepsilon}$ tales que
\[
m(E \setminus F_{m}) < \frac{\varepsilon}{2^{m}}
\]
y
\[
|f(x) - f_{k}(x)| < \frac{1}{m} \quad \text{para todo } k \geq \kappa_{m}^{\varepsilon} \text{ y } x \in F_{m}.
\]
Sea
\[
F_{\varepsilon} = \bigcap_{m=1}^{\infty} F_{m}.
\]
Entonces $F_{\varepsilon}$ es cerrado, por ser intersección de cerrados, y
\[
m(E \setminus F_{\varepsilon}) = m\left( E \cap \bigcup_{m=1}^{\infty} F_{m}^{c} \right) \leq \sum_{m=1}^{\infty} m(E \setminus F_{m}) < \sum_{m=1}^{\infty} \frac{\varepsilon}{2^{m}} = \varepsilon.
\]
Veamos que la convergencia es uniforme en $F_{\varepsilon}$. Dado $\delta > 0$, tome $m \geq 1$ con $\tfrac{1}{m} < \delta$. Si $x \in F_{\varepsilon} \subseteq F_{m}$ y $k \geq \kappa_{m}^{\varepsilon}$, entonces
\[
|f(x) - f_{k}(x)| < \frac{1}{m} < \delta.
\]
Como $\kappa_{m}^{\varepsilon}$ no depende de $x$, concluimos que $f_{k} \to f$ uniformemente en $F_{\varepsilon}$.
\end{proof}

\subsection{El teorema de Lusin}

\begin{namedstd}{Definición}{Propiedad $C$ de una función}
Una función $f : E \to \R$ tiene la \emph{propiedad $C$} en $E$ si dado $\varepsilon > 0$, existe un $F \subseteq E$ cerrado tal que
\begin{enumerate}[label=(\arabic*)]
\item $m(E \setminus F) < \varepsilon$;
\item $f$ es continua relativa a $F$, es decir, $f : F \to \R$ es continua.
\end{enumerate}
\end{namedstd}

\begin{namedstd}{Nota}{Caracterización secuencial de la continuidad relativa}
La segunda condición de la definición anterior es equivalente a que si $\{x_{n}\}_{n=1}^{\infty} \subseteq F$ con
\[
\lim_{n \to \infty} x_{n} = x \in F,
\]
entonces
\[
\lim_{n \to \infty} f(x_{n}) = f(x).
\]
\end{namedstd}

\begin{namedres}{Lema}{Las funciones simples medibles tienen la propiedad $C$}
Sea $\phi$ una función simple y medible en $E$. Entonces $\phi$ tiene la propiedad $C$.
\end{namedres}
\begin{proof}
Escribamos
\[
\phi(x) = \sum_{\ell=1}^{m} b_{\ell} \ind_{B_{\ell}}(x)
\]
con los $B_{\ell} \subseteq E$ medibles, $B_{i} \cap B_{j} = \emptyset$ y $b_{i} \neq b_{j}$ si $i \neq j$, y $E = \bigcup_{\ell=1}^{m} B_{\ell}$.

Dado $\varepsilon > 0$, por regularidad interna tome, para cada $1 \leq \ell \leq m$, un cerrado $F_{\ell} \subseteq B_{\ell}$ tal que
\[
m(B_{\ell} \setminus F_{\ell}) < \frac{\varepsilon}{m}.
\]
Entonces
\[
F = \bigcup_{\ell=1}^{m} F_{\ell}
\]
es cerrado, por ser unión finita de cerrados. Además, como los $B_{\ell}$ cubren $E$, tenemos que $E \setminus F \subseteq \bigcup_{\ell=1}^{m} (B_{\ell} \setminus F_{\ell})$, y por lo tanto
\[
m(E \setminus F) \leq \sum_{\ell=1}^{m} m(B_{\ell} \setminus F_{\ell}) < m \cdot \frac{\varepsilon}{m} = \varepsilon.
\]
Falta ver que $\phi$ es continua relativa a $F$. Tome $\{x_{n}\}_{n=1}^{\infty} \subseteq F$ tal que
\[
\lim_{n \to \infty} x_{n} = y \in F.
\]
Entonces existe $\ell_{0}$ tal que $y \in F_{\ell_{0}}$. Afirmamos que, dado $\ell \neq \ell_{0}$, el conjunto
\[
\{x_{n}\}_{n=1}^{\infty} \cap F_{\ell}
\]
es finito. En caso contrario, existe una subsucesión
\[
\{x_{n_{k}}\}_{k=1}^{\infty} \subseteq F_{\ell} \quad \text{con} \quad \lim_{k \to \infty} x_{n_{k}} = y.
\]
Como $F_{\ell}$ es cerrado, esto implica que $y \in F_{\ell}$; pero $F_{\ell} \cap F_{\ell_{0}} \subseteq B_{\ell} \cap B_{\ell_{0}} = \emptyset$, lo que nos lleva a una contradicción.

Por lo tanto, como los $x_{n}$ están repartidos entre los $F_{\ell}$ y solo finitos caen fuera de $F_{\ell_{0}}$, existe $k_{0}$ tal que $x_{n} \in F_{\ell_{0}}$ para $n \geq k_{0}$. Entonces
\[
\phi(x_{n}) = b_{\ell_{0}} = \phi(y) \quad \text{para } n \geq k_{0},
\]
es decir,
\[
\lim_{n \to \infty} \phi(x_{n}) = \phi(y). \qedhere
\]
\end{proof}

\begin{namedres}{Teorema}{Lusin: medibilidad equivale a la propiedad $C$}
Sea $f : E \to \R$ con $E$ medible. Entonces $f$ es medible si y solo si $f$ tiene la propiedad $C$ en $E$.
\end{namedres}
\begin{proof}
\textit{Medible implica propiedad $C$, caso $m(E) < \infty$.} Sea $f : E \to \R$ medible; entonces existe $\{f_{k}\}_{k=1}^{\infty}$ una sucesión de funciones simples y medibles tales que
\[
\lim_{k \to \infty} f_{k}(x) = f(x)
\]
c.p.d. en $E$. Sea $\varepsilon > 0$. Por el lema anterior, para cada $k \geq 1$ existe $F_{k} \subseteq E$ cerrado que satisface
\[
m(E \setminus F_{k}) < \frac{\varepsilon}{2^{k+1}},
\]
y tal que $f_{k}$ es continua relativa a $F_{k}$. Como $m(E) < \infty$, por el teorema de Egorov existe $F_{0} \subseteq E$ cerrado tal que
\[
m(E \setminus F_{0}) < \frac{\varepsilon}{2} \qquad \text{y} \qquad \lim_{k \to \infty} f_{k} = f \ \text{uniformemente en } F_{0}.
\]
Si tomamos
\[
F = F_{0} \cap \bigcap_{k=1}^{\infty} F_{k},
\]
entonces $F$ es cerrado, y como
\[
E \setminus F = (E \setminus F_{0}) \cup \bigcup_{k=1}^{\infty} (E \setminus F_{k}),
\]
tenemos que
\[
m(E \setminus F) \leq m(E \setminus F_{0}) + \sum_{k=1}^{\infty} m(E \setminus F_{k}) < \frac{\varepsilon}{2} + \sum_{k=1}^{\infty} \frac{\varepsilon}{2^{k+1}} = \varepsilon.
\]
Cada $f_{k}$ es continua relativa a $F_{k} \supseteq F$, y por lo tanto continua relativa a $F$. Como
\[
\lim_{k \to \infty} f_{k} = f
\]
uniformemente en $F$, y el límite uniforme de funciones continuas es continuo, tenemos que $f$ es continua relativa a $F$.

\textit{Medible implica propiedad $C$, caso $m(E) = \infty$.} Llamemos
\[
E_{k} = E \cap \{ x \in \R^{d} :\ k-1 \leq |x| < k \}, \quad k \geq 1.
\]
Los $E_{k}$ son medibles, disjuntos dos a dos, cubren $E$ y $m(E_{k}) \leq m(B(0,k)) < \infty$. Por el caso anterior, para cada $k$ existe $F_{k} \subseteq E_{k}$ cerrado tal que $f$ es continua relativa a $F_{k}$ y
\[
m(E_{k} \setminus F_{k}) \leq \frac{\varepsilon}{2^{k}}.
\]
Tome
\[
F = \bigcup_{k=1}^{\infty} F_{k}.
\]
Como $E \setminus F \subseteq \bigcup_{k=1}^{\infty} (E_{k} \setminus F_{k})$, tenemos que $m(E \setminus F) \leq \sum_{k=1}^{\infty} \varepsilon/2^{k} = \varepsilon$.

Veamos que $F$ es cerrado. Sea $\{x_{n}\}_{n=1}^{\infty} \subseteq F$ con $\lim_{n \to \infty} x_{n} = y \in \R^{d}$. Como la sucesión converge, es acotada: existe $k_{0}$ tal que
\[
n \geq k_{0} \implies |x_{n}| \leq |y| + 1.
\]
Sea $M$ un entero con $M \geq |y| + 2$. Si $x_{n} \in F_{k} \subseteq E_{k}$ con $n \geq k_{0}$, entonces $k - 1 \leq |x_{n}| \leq |y| + 1$, es decir $k \leq |y| + 2 \leq M$. Por lo tanto
\[
x_{n} \in \bigcup_{k=1}^{M} F_{k} \quad \text{para } n \geq k_{0},
\]
que es una unión finita de cerrados y por ende un cerrado. Entonces $y \in \bigcup_{k=1}^{M} F_{k} \subseteq F$, y $F$ es cerrado.

Veamos que $f$ es continua relativa a $F$. Sea $\{x_{n}\}_{n=1}^{\infty} \subseteq F$ con $\lim_{n \to \infty} x_{n} = y \in F$. Por lo anterior, existe $k_{0}$ tal que $x_{n} \in \bigcup_{k=1}^{M} F_{k}$ para $n \geq k_{0}$. Como los $F_{k}$ son cerrados y disjuntos dos a dos, por el argumento de la prueba del lema anterior existen $k_{1}$ y $\ell_{1}$ tales que
\[
n \geq k_{1} \implies x_{n} \in F_{\ell_{1}},
\]
con $y \in F_{\ell_{1}}$. Como $f$ es continua relativa a $F_{\ell_{1}}$, concluimos que $\lim_{n \to \infty} f(x_{n}) = f(y)$.

\textit{Propiedad $C$ implica medible.} Si $f$ posee la propiedad $C$, entonces para cada $k \geq 1$ existe $F_{k} \subseteq E$ cerrado que satisface:
\begin{enumerate}[label=(\roman*)]
\item $m(E \setminus F_{k}) < \frac{1}{k}$;
\item $f$ es continua relativa a $F_{k}$.
\end{enumerate}
Sea
\[
H = \bigcup_{k=1}^{\infty} F_{k},
\]
entonces $H \subseteq E$ y $m(E \setminus H) = 0$ (ejercicio; véase el ejercicio siguiente). Escribamos $Z = E \setminus H$. Finalmente, para cada $a \in \R$,
\begin{align*}
\{ x \in E :\ f(x) > a \} &= \{ x \in H :\ f(x) > a \} \cup \{ x \in Z :\ f(x) > a \} \\
&= \bigcup_{k=1}^{\infty} \{ x \in F_{k} :\ f(x) > a \} \cup \{ x \in Z :\ f(x) > a \}.
\end{align*}
Cada $\{ x \in F_{k} :\ f(x) > a \}$ es medible: como $f$ es continua relativa a $F_{k}$, este conjunto es abierto relativo a $F_{k}$, es decir, de la forma $F_{k} \cap G$ con $G$ abierto de $\R^{d}$, que es medible. Por otra parte, $\{ x \in Z :\ f(x) > a \} \subseteq Z$ tiene medida exterior cero y es por lo tanto medible. Entonces $\{ f > a \}$ es medible para todo $a \in \R$, es decir, $f$ es medible.
\end{proof}

\begin{namedstd}{Ejercicio}{El complemento de la unión de los $F_{k}$ es nulo}
Con la notación de la prueba anterior, muestre que si $H = \bigcup_{k=1}^{\infty} F_{k}$ con $m(E \setminus F_{k}) < \tfrac{1}{k}$ para todo $k \geq 1$, entonces $m(E \setminus H) = 0$.
\end{namedstd}

% =============================================================================
\section{Convergencia en medida}
% =============================================================================

\subsection{Definición y relación con la convergencia c.p.d.}

\begin{namedstd}{Definición}{Convergencia en medida}
Sean $f, f_{n} : E \to \overline{\R}$ medibles y finitas casi por doquier. Decimos que $f_{n}$ \emph{converge a $f$ en medida} si, para todo $\varepsilon > 0$,
\[
\lim_{n \to \infty} m\{ |f_{n} - f| > \varepsilon \} = 0.
\]
Denotamos $f_{n} \xrightarrow{m} f$.
\end{namedstd}

\begin{namedstd}{Nota}{La convergencia en medida es la más débil del curso}
Esta convergencia es la convergencia más débil que veremos en el curso: la convergencia uniforme implica la convergencia puntual, la puntual implica la convergencia c.p.d., y sobre conjuntos de medida finita la convergencia c.p.d. implica la convergencia en medida, como muestra el teorema siguiente. Ninguna de las implicaciones recíprocas vale en general.
\end{namedstd}

\begin{namedres}{Teorema}{En medida finita, la convergencia c.p.d. implica convergencia en medida}
Sean $f, f_{n} : E \to \overline{\R}$ medibles y finitas casi por doquier. Si
\[
\lim_{n \to \infty} f_{n} = f
\]
casi por doquier en $E$ y $m(E) < \infty$, entonces $f_{n} \xrightarrow{m} f$ en $E$.
\end{namedres}
\begin{proof}
Sean $\eta > 0$ y $\varepsilon > 0$; hay que mostrar que existe $n_{0}$ tal que
\[
m\{ |f_{n} - f| > \varepsilon \} < \eta
\]
cuando $n \geq n_{0}$. Como $m(E) < \infty$, por el teorema de Egorov existe $F \subseteq E$ cerrado tal que
\[
m(E \setminus F) < \eta
\]
y $f_{n} \to f$ uniformemente en $F$. Entonces existe $n_{0}$ tal que
\[
|f_{n}(x) - f(x)| < \varepsilon \quad \text{para todo } n \geq n_{0} \text{ y } x \in F.
\]
Luego
\[
\{ x \in E :\ |f_{n}(x) - f(x)| > \varepsilon \} \subseteq E \setminus F
\]
cuando $n \geq n_{0}$, y por monotonía de la medida,
\[
m\{ x \in E :\ |f_{n}(x) - f(x)| > \varepsilon \} \leq m(E \setminus F) < \eta. \qedhere
\]
\end{proof}

\begin{namedstd}{Ejemplo}{La máquina de escribir: convergencia en medida sin convergencia puntual}
Considere la sucesión de indicadoras de intervalos diádicos de $[0,1]$:
\begin{gather*}
I_{0} = \ind_{[0,1]}, \\
I_{1,1} = \ind_{[0,\frac{1}{2}]}, \quad I_{1,2} = \ind_{[\frac{1}{2},1]}, \\
I_{2,i} = \ind_{[\frac{i-1}{2^{2}},\frac{i}{2^{2}}]}, \quad 1 \leq i \leq 4.
\end{gather*}
En general, $I_{n,i} = \ind_{[\frac{i-1}{2^{n}},\frac{i}{2^{n}}]}$ con $1 \leq i \leq 2^{n}$. Para obtener una sucesión ordenamos los índices lexicográficamente:
\[
(m,i) \leq (n,j) \iff (m < n) \lor \big( (m = n) \land i \leq j \big).
\]
Esta sucesión converge a cero en medida, pues para $0 < \varepsilon < 1$ el conjunto $\{ I_{n,i} > \varepsilon \}$ es un intervalo de medida $2^{-n}$, y $n \to \infty$ a lo largo de la sucesión. Sin embargo, no converge en ningún punto: para cada $x \in [0,1]$ y cada $n$ existe algún $i$ con $x \in [\frac{i-1}{2^{n}},\frac{i}{2^{n}}]$, de modo que la sucesión toma los valores $1$ y $0$ infinitas veces en $x$.
\end{namedstd}

\begin{namedres}{Teorema}{Toda sucesión convergente en medida tiene una subsucesión convergente c.p.d.}
Sea $f_{n} \xrightarrow{m} f$ en $E$. Entonces existen índices $n_{1} < n_{2} < \dots$ de manera que
\[
f_{n_{j}} \to f
\]
casi por doquier en $E$.
\end{namedres}
\begin{proof}
Dado $j \geq 1$, por la convergencia en medida existe $n_{j}$ tal que
\[
m\left\{ |f_{n} - f| > \frac{1}{j} \right\} \leq \frac{1}{2^{j}} \quad \text{si } n \geq n_{j}.
\]
Sin pérdida de generalidad podemos suponer $n_{j} < n_{j+1}$, escogiendo los índices de manera creciente. Entonces
\[
m\left\{ |f_{n_{j}} - f| > \frac{1}{j} \right\} \leq \frac{1}{2^{j}}.
\]
Tome
\[
H_{m} = \bigcup_{j \geq m} \left\{ |f_{n_{j}} - f| > \frac{1}{j} \right\}.
\]
Entonces $H_{m+1} \subseteq H_{m}$ y, por subaditividad numerable,
\[
m(H_{m}) \leq \sum_{j \geq m} m\left\{ |f_{n_{j}} - f| > \frac{1}{j} \right\} \leq \sum_{j \geq m} \frac{1}{2^{j}} = \frac{1}{2^{m-1}}.
\]
Considere
\[
Z = \bigcap_{m \in \N} H_{m}.
\]
Como $Z \subseteq H_{m}$ para todo $m$, tenemos que $m(Z) \leq 2^{-(m-1)}$ para todo $m$, y entonces $Z$ tiene medida cero. Además, si
\[
x \in E \setminus Z = \bigcup_{m \in \N} (E \setminus H_{m}),
\]
entonces existe $m_{0}$ tal que $x \in E \setminus H_{m_{0}}$. Como
\[
E \setminus H_{m_{0}} = E \setminus \bigcup_{j \geq m_{0}} \left\{ |f_{n_{j}} - f| > \frac{1}{j} \right\} = \bigcap_{j \geq m_{0}} \left\{ |f_{n_{j}} - f| \leq \frac{1}{j} \right\},
\]
obtenemos que, dado $x \in E \setminus Z$, existe $m_{0}$ tal que
\[
|f_{n_{j}}(x) - f(x)| \leq \frac{1}{j} \quad \text{para } j \geq m_{0}.
\]
Finalmente, dado $\varepsilon > 0$, existe $j_{0}$ tal que $\tfrac{1}{j_{0}} < \varepsilon$, y entonces
\[
|f_{n_{j}}(x) - f(x)| \leq \frac{1}{j} \leq \frac{1}{j_{0}} < \varepsilon
\]
para $j \geq \max(j_{0}, m_{0})$. Es decir, $f_{n_{j}}(x) \to f(x)$ para todo $x$ fuera del conjunto nulo $Z$.
\end{proof}

\subsection{Sucesiones de Cauchy en medida}

\begin{namedstd}{Definición}{Sucesión de Cauchy en medida}
Decimos que $(f_{n})_{n \in \N}$ es \emph{de Cauchy en medida} si para todos $\varepsilon, \eta > 0$ existe $N \in \N$ tal que
\[
m\{ |f_{n} - f_{m}| > \varepsilon \} < \eta
\]
para $n, m \geq N$.
\end{namedstd}

\begin{namedres}{Lema}{Toda sucesión de Cauchy en medida converge en medida}
Una sucesión de Cauchy en medida es convergente en medida.
\end{namedres}
\begin{proof}
Ejercicio.
\end{proof}

% =============================================================================
\section{La integral de Lebesgue de funciones no negativas}
% =============================================================================

\subsection{La región bajo el gráfico}

\begin{namedstd}{Definición}{Región bajo el gráfico y gráfico de una función}
Sea $f : E \to \overline{\R}$ medible con $E \subseteq \R^{d}$ medible y $f \geq 0$. Definimos la \emph{región bajo el gráfico} de $f$ sobre $E$ como
\[
R(f,E) = \{ (x,y) \in \R^{d+1} :\ x \in E,\ 0 \leq y \leq f(x) \},
\]
y el \emph{gráfico} de $f$ sobre $E$ como
\[
\Gamma(f,E) = \{ (x, f(x)) :\ x \in E \}.
\]
Surge la pregunta: ¿es $R(f,E)$ medible?
\end{namedstd}

\begin{namedstd}{Ejemplo}{La región bajo un múltiplo de una indicadora}
Analicemos el caso
\[
f(x) = a \ind_{A}(x)
\]
con $A \subseteq E$ medible y $a > 0$. Entonces ocurre que
\begin{align*}
R(f,E) &= \{ (x,y) \in \R^{d+1} :\ x \in E,\ 0 \leq y \leq f(x) \} \\
&= \{ (x,y) \in \R^{d+1} :\ x \in A,\ 0 \leq y \leq a \} \cup \{ (x,y) \in \R^{d+1} :\ x \in E \setminus A,\ y = 0 \} \\
&= \big( A \times [0,a] \big) \cup \big( (E \setminus A) \times \{0\} \big).
\end{align*}
Es decir, la medibilidad de $R(f,E)$ en este caso se reduce a la medibilidad de productos de la forma $A \times [0,a]$.
\end{namedstd}

\begin{namedres}{Lema}{Medida del producto de un medible por un intervalo}
Sea $A \subseteq \R^{d}$ medible y $a \geq 0$. Entonces $A \times [0,a] \subseteq \R^{d+1}$ es medible y su medida es
\[
m(A \times [0,a]) = a\, m(A),
\]
con la convención $0 \cdot \infty = 0$.
\end{namedres}
\begin{proof}
\textit{Caso de una caja.} Si $A = [a_{1},b_{1}] \times \dots \times [a_{d},b_{d}]$, entonces $A \times [0,a]$ es una caja de $\R^{d+1}$ y el resultado se sigue de la fórmula para el volumen de cajas. Verificar los detalles queda como ejercicio para la persona lectora.

\textit{Caso de un abierto con $m(A) < \infty$.} Si $A$ es abierto, entonces existen cajas $I_{k}$ en $d$ dimensiones tales que
\[
A = \bigcup_{k=1}^{\infty} I_{k} \quad \text{con } I_{j}^{\circ} \cap I_{k}^{\circ} = \emptyset \text{ si } j \neq k.
\]
Entonces
\[
A \times [0,a] = \bigcup_{k=1}^{\infty} I_{k} \times [0,a]
\]
es un conjunto medible, por ser unión numerable de cajas. Como
\[
(I_{k} \times [0,a])^{\circ} \cap (I_{j} \times [0,a])^{\circ} = \emptyset, \quad k \neq j,
\]
y la medida de una unión numerable de cajas con interiores disjuntos dos a dos es la suma de sus medidas, tenemos que
\[
m(A \times [0,a]) = \sum_{i=1}^{\infty} m(I_{i} \times [0,a]) = a \sum_{i=1}^{\infty} m(I_{i}) = a\, m(A).
\]

\textit{Caso de un $G_{\delta}$ con $m(A) < \infty$.} Sea ahora
\[
A = \bigcap_{j=1}^{\infty} G_{j}
\]
un $G_{\delta}$ con cada $G_{i}$ abierto y $G_{i+1} \subseteq G_{i}$. Podemos asumir la monotonía porque, reemplazando $G_{j}$ por $G_{1} \cap \dots \cap G_{j}$, que sigue siendo abierto, la intersección total no cambia. Además, como $m(A) < \infty$, por regularidad exterior existe un abierto $G \supseteq A$ con $m(G) < \infty$; intersecando cada $G_{j}$ con $G$ podemos suponer también que $m(G_{1}) < \infty$. Luego
\[
A \times [0,a] = \bigcap_{j=1}^{\infty} G_{j} \times [0,a],
\]
con $G_{i+1} \times [0,a] \subseteq G_{i} \times [0,a]$ medibles y de medida finita, pues $m(G_{i} \times [0,a]) = a\, m(G_{i}) \leq a\, m(G_{1}) < \infty$ por el caso anterior. Entonces, por continuidad desde arriba de la medida,
\[
m(A \times [0,a]) = \lim_{i \to \infty} m(G_{i} \times [0,a]) = \lim_{i \to \infty} a\, m(G_{i}) = a\, m(A).
\]

\textit{Caso medible general con $m(A) < \infty$.} En el caso de que $A$ sea medible con medida finita, podemos escribir
\[
A = H \setminus Z
\]
con $H$ un conjunto $G_{\delta}$ tal que $A \subseteq H$, $m(H) = m(A)$, y $Z = H \setminus A$ de medida cero. Entonces, por el paso anterior, $H \times [0,a]$ es medible y
\[
m(H \times [0,a]) = a\, m(H).
\]
Se deja de ejercicio probar que
\[
m_{e}(Z \times [0,a]) = 0.
\]
Entonces $Z \times [0,a]$ es medible con medida cero,
\[
A \times [0,a] = (H \times [0,a]) \setminus (Z \times [0,a])
\]
es medible, y
\[
m(A \times [0,a]) = m(H \times [0,a]) = a\, m(H) = a\, m(A).
\]

\textit{Caso $m(A) = \infty$.} Finalmente, si $m(A) = \infty$, llamemos
\[
A_{k} = A \cap B(0,k).
\]
De esta manera
\[
A = \bigcup_{k=1}^{\infty} A_{k},
\]
con $A_{k} \subseteq A_{k+1}$ medibles y acotados, en particular de medida finita. Por los argumentos anteriores,
\[
A_{k} \times [0,a] \subseteq A_{k+1} \times [0,a]
\]
son conjuntos medibles tales que
\[
m(A_{k} \times [0,a]) = a\, m(A_{k}).
\]
Entonces $A \times [0,a] = \bigcup_{k=1}^{\infty} A_{k} \times [0,a]$ es medible y, por continuidad desde abajo,
\[
m(A \times [0,a]) = \lim_{k \to \infty} m(A_{k} \times [0,a]) = \lim_{k \to \infty} a\, m(A_{k}) = a\, m(A). \qedhere
\]
\end{proof}

\begin{namedstd}{Ejercicio}{El producto de un nulo por un intervalo es nulo}
Sea $Z \subseteq \R^{d}$ con $m(Z) = 0$ y $a \geq 0$. Muestre que $m_{e}(Z \times [0,a]) = 0$.
\end{namedstd}

\subsection{El gráfico de una función medible}

\begin{namedres}{Lema}{El gráfico de una función medible tiene medida cero}
Sea $f : E \to \R$ medible, con $E$ medible. Entonces $m_{e}(\Gamma(f,E)) = 0$.
\end{namedres}
\begin{proof}
Asumamos primero que $E$ tiene medida finita. Sea $\varepsilon > 0$ y, para cada $k \in \Z$,
\[
E_{k} = \{ x \in E :\ k \varepsilon \leq f(x) < (k+1) \varepsilon \}.
\]
Los $E_{k}$ son medibles, disjuntos dos a dos, y
\[
E = \bigcup_{k \in \Z} E_{k}.
\]
Note que
\[
\Gamma(f, E_{k}) = \{ (x, f(x)) :\ x \in E_{k} \} \subseteq E_{k} \times [k\varepsilon, (k+1)\varepsilon).
\]
Por el lema anterior y la invariancia de la medida exterior bajo traslaciones, tenemos que
\[
m_{e}(\Gamma(f,E_{k})) \leq m_{e}\big( E_{k} \times [k\varepsilon,(k+1)\varepsilon) \big) \leq \varepsilon\, m(E_{k}).
\]
Por lo tanto, por subaditividad numerable de la medida exterior, vale que
\[
m_{e}(\Gamma(f,E)) = m_{e}\left( \bigcup_{k \in \Z} \Gamma(f,E_{k}) \right) \leq \sum_{k \in \Z} m_{e}(\Gamma(f,E_{k})) \leq \sum_{k \in \Z} \varepsilon\, m(E_{k}) = \varepsilon\, m(E),
\]
donde la última igualdad usa que los $E_{k}$ son disjuntos dos a dos. Como $\varepsilon > 0$ es arbitrario y $m(E) < \infty$, concluimos que $m_{e}(\Gamma(f,E)) = 0$.

Terminar la prueba en el caso $m(E) = \infty$ queda como ejercicio (véase el ejercicio siguiente).
\end{proof}

\begin{namedstd}{Ejercicio}{El caso de medida infinita del lema del gráfico}
Termine la prueba del lema anterior: muestre que si $m(E) = \infty$, entonces también $m_{e}(\Gamma(f,E)) = 0$.
\end{namedstd}

\begin{namedres}{Lema}{La región bajo una función simple no negativa es medible}
Sea $\phi : E \to \R$ simple, medible y no negativa, digamos
\[
\phi(x) = \sum_{i=1}^{m} a_{i} \ind_{A_{i}},
\]
con $a_{i} \neq a_{j}$ y $A_{i} \cap A_{j} = \emptyset$ para $i \neq j$, y $E = \bigcup_{i=1}^{m} A_{i}$. Entonces $R(\phi,E)$ es medible.
\end{namedres}
\begin{proof}
Tenemos que
\[
R(\phi,E) = \{ (x,y) \in \R^{d+1} :\ x \in E,\ 0 \leq y \leq \phi(x) \} = \bigcup_{j=1}^{m} \{ (x,y) \in \R^{d+1} :\ x \in A_{j},\ 0 \leq y \leq a_{j} \},
\]
es decir, $R(\phi,E) = \bigcup_{j=1}^{m} A_{j} \times [0,a_{j}]$, que es una unión finita de conjuntos medibles por el lema del producto, y por lo tanto es un conjunto medible.
\end{proof}

\begin{namedres}{Teorema}{La región bajo el gráfico de una función medible no negativa es medible}
Sea $f : E \to \R$ medible con $f \geq 0$ y $E$ medible. Entonces $R(f,E) \subseteq \R^{d+1}$ es un conjunto medible.
\end{namedres}
\begin{proof}
Como $f \geq 0$ es medible, existe una sucesión creciente de funciones simples, medibles y no negativas $\{\phi_{k}\}_{k=1}^{\infty}$ que satisface
\[
\lim_{k \to \infty} \phi_{k}(x) = f(x)
\]
en $E$. Luego, si $0 \leq y < f(x)$, como $\phi_{k}(x) \to f(x) > y$, existe $\phi_{k}$ tal que
\[
0 \leq y \leq \phi_{k}(x).
\]
Entonces, separando los puntos con $y < f(x)$ de aquellos con $y = f(x)$,
\[
R(f,E) = \{ (x,y) :\ x \in E,\ 0 \leq y < f(x) \} \cup \{ (x,y) :\ x \in E,\ f(x) = y \},
\]
y por lo anterior obtenemos
\[
R(f,E) = \bigcup_{k=1}^{\infty} R(\phi_{k},E) \cup \{ (x,y) :\ x \in E,\ f(x) = y,\ y \geq 0 \},
\]
pues cada $R(\phi_{k},E) \subseteq R(f,E)$ al ser $\phi_{k} \leq f$. Ahora, cada $R(\phi_{k},E)$ es medible por el lema anterior, y el segundo conjunto está contenido en $\Gamma(f,E)$, que tiene medida exterior cero y es por lo tanto medible. Concluimos que $R(f,E)$ es medible, por ser unión numerable de medibles.
\end{proof}

\subsection{La definición de la integral}

\begin{namedstd}{Definición}{Integral de Lebesgue de una función medible no negativa}
Dada $f : E \to \R$ medible tal que $f \geq 0$, definimos su \emph{integral de Lebesgue} como
\[
\int_{E} f \, dx = m(R(f,E)).
\]
\end{namedstd}

\begin{namedres}{Proposición}{La integral de una función simple no negativa}
Sea $\phi = \sum_{i=1}^{m} a_{i} \ind_{A_{i}}$ simple, medible y no negativa, con los $A_{i} \subseteq E$ medibles y disjuntos dos a dos. Entonces
\[
\int_{E} \phi(x) \, dx = \sum_{i=1}^{m} a_{i}\, m(A_{i}).
\]
\end{namedres}
\begin{proof}
Podemos asumir que $a_{i} \neq 0$ para $1 \leq i \leq m$, pues los términos con $a_{i} = 0$ no aportan a ninguno de los dos lados. Entonces
\begin{align*}
\int_{E} \phi(x)\, dx &= m\big( R(\phi,E) \big) \\
&= m\big( \{ \phi = 0 \} \times \{0\} \big) + m\left( \bigcup_{i=1}^{m} \{ (x,y) :\ x \in A_{i},\ 0 \leq y \leq a_{i} \} \right) \\
&= 0 + \sum_{i=1}^{m} a_{i}\, m(A_{i}),
\end{align*}
donde usamos que $\{\phi = 0\} \times \{0\}$ tiene medida cero, que los conjuntos $A_{i} \times [0,a_{i}]$ son disjuntos dos a dos porque los $A_{i}$ lo son, y la fórmula $m(A_{i} \times [0,a_{i}]) = a_{i}\, m(A_{i})$.
\end{proof}

\begin{namedres}{Teorema}{Monotonía de la integral respecto al integrando y al dominio}
Sean $f, g : E \to [0,\infty)$ medibles.
\begin{enumerate}[label=(\roman*)]
\item Si $0 \leq g \leq f$, entonces
\[
\int_{E} g(x) \, dx \leq \int_{E} f(x) \, dx.
\]
\item Si $E_{1} \subseteq E_{2} \subseteq E$ son medibles, entonces
\[
\int_{E_{1}} f(x) \, dx \leq \int_{E_{2}} f(x) \, dx.
\]
\end{enumerate}
\end{namedres}
\begin{proof}
Dado que $0 \leq g \leq f$, tenemos que
\[
\{ (x,y) :\ x \in E,\ 0 \leq y \leq g(x) \} \subseteq \{ (x,y) :\ x \in E,\ 0 \leq y \leq f(x) \},
\]
luego $R(g,E) \subseteq R(f,E)$, y la parte (i) se sigue de la monotonía de la medida.

Para la segunda parte vamos primero a probar que
\[
\int_{E_{1}} f(x) \, dx = \int_{E} \ind_{E_{1}} f(x) \, dx
\]
para $E_{1} \subseteq E$ medible. Note que
\[
\{ (x,y) :\ x \in E,\ 0 \leq y \leq \ind_{E_{1}}(x) f(x) \} = \big( (E \setminus E_{1}) \times \{0\} \big) \cup \{ (x,y) :\ x \in E_{1},\ 0 \leq y \leq f(x) \},
\]
y como $(E \setminus E_{1}) \times \{0\}$ tiene medida cero, entonces
\[
m\big( R(\ind_{E_{1}} f, E) \big) = m\big( R(f,E_{1}) \big).
\]
Ahora, si $E_{1} \subseteq E_{2}$ son medibles, entonces
\[
\ind_{E_{1}} f \leq \ind_{E_{2}} f,
\]
y por la parte (i) obtenemos
\[
\int_{E} \ind_{E_{1}} f(x) \, dx \leq \int_{E} \ind_{E_{2}} f(x) \, dx.
\]
Concluimos que
\[
\int_{E_{1}} f(x) \, dx \leq \int_{E_{2}} f(x) \, dx. \qedhere
\]
\end{proof}

% =============================================================================
\section{Propiedades de la integral de funciones no negativas}
% =============================================================================

\subsection{Funciones con valores infinitos}

\begin{namedstd}{Definición}{Integral de una función medible con valores en $[0,\infty]$}
Sea $f : E \to [0,\infty]$ medible con $E$ medible y $f \geq 0$. Tome
\[
E_{1} = \{ x \in E :\ f(x) = \infty \}.
\]
Definimos
\begin{align*}
R(f,E) &= \{ (x,y) \in \R^{d+1} :\ x \in E,\ 0 \leq y \leq f(x) \} \\
&= \{ (x,y) \in \R^{d+1} :\ x \in E \setminus E_{1},\ 0 \leq y \leq f(x) \} \cup \big( E_{1} \times [0,\infty) \big),
\end{align*}
y
\[
\Gamma(f,E) = \{ (x,f(x)) \in \R^{d+1} :\ x \in E \setminus E_{1} \}.
\]
Al igual que se hizo para funciones de valores reales, definimos
\[
\int_{E} f \, dx = m\big( R(f,E) \big).
\]
\end{namedstd}

\begin{namedstd}{Nota}{La región sigue siendo medible para valores infinitos}
La teoría desarrollada hasta ahora asegura que $R(f,E)$ es medible siempre que $f$ y $E$ sean medibles: la primera parte de la descomposición es la región bajo la restricción de $f$ al medible $E \setminus E_{1}$, y
\[
E_{1} \times [0,\infty) = \bigcup_{n=1}^{\infty} E_{1} \times [0,n]
\]
es una unión numerable de conjuntos medibles.
\end{namedstd}

\begin{namedres}{Proposición}{Una función con integral finita es finita c.p.d.}
Sea $f : E \to [0,\infty]$ medible tal que
\[
\int_{E} f(x) \, dx < \infty.
\]
Entonces $m(\{ x \in E : f(x) = \infty \}) = 0$, es decir, $f$ es finita casi por doquier.
\end{namedres}
\begin{proof}
Sea $E_{1} = \{ f = \infty \}$ y, para cada $n \in \N$, tome
\[
g(x) = n \ind_{E_{1}}(x).
\]
Entonces $g(x) \leq f(x)$ para $x \in E$, pues en $E_{1}$ se tiene $f = \infty \geq n$ y fuera de $E_{1}$ se tiene $g = 0 \leq f$. Luego, por la fórmula para funciones simples y la monotonía de la integral,
\[
n\, m(E_{1}) = \int_{E} g(x)\, dx \leq \int_{E} f(x) \, dx < \infty
\]
para todo $n \in \N$, y por lo tanto $m(E_{1}) = 0$.
\end{proof}

\subsection{Descomposición del dominio y el teorema de convergencia monótona}

\begin{namedres}{Proposición}{Aditividad de la integral respecto al dominio}
Sea $f : E \to [0,\infty]$ medible. En el caso de que $E = \bigcup_{i=1}^{\infty} E_{i}$ con los $E_{i}$ medibles y $E_{i} \cap E_{j} = \emptyset$ si $i \neq j$, tenemos que
\[
\int_{E} f(x) \, dx = \sum_{k=1}^{\infty} \int_{E_{k}} f(x) \, dx.
\]
Esta fórmula es válida también para uniones finitas.
\end{namedres}
\begin{proof}
Como los $E_{i}$ son disjuntos dos a dos, tenemos que
\[
R(f,E) = \bigcup_{k=1}^{\infty} R(f,E_{k}),
\]
y los $R(f,E_{k})$ son disjuntos dos a dos, pues sus primeras coordenadas viven en conjuntos disjuntos. De esta manera, por aditividad numerable de la medida,
\[
\int_{E} f(x) \, dx = m(R(f,E)) = m\left( \bigcup_{k=1}^{\infty} R(f,E_{k}) \right) = \sum_{k=1}^{\infty} m(R(f,E_{k})) = \sum_{k=1}^{\infty} \int_{E_{k}} f(x) \, dx.
\]
La fórmula para uniones finitas se obtiene tomando $E_{k} = \emptyset$ para los índices restantes, pues $R(f,\emptyset) = \emptyset$.
\end{proof}

\begin{namedres}{Teorema}{Convergencia monótona para funciones no negativas}
Sea $\{f_{k}\}_{k=1}^{\infty}$ una sucesión de funciones medibles tales que
\[
0 \leq f_{k}(x) \leq f_{k+1}(x)
\]
para todo $x \in E$. Si
\[
\lim_{k \to \infty} f_{k}(x) = f(x)
\]
casi por doquier en $E$, entonces
\[
\lim_{k \to \infty} \int_{E} f_{k}(x) \, dx = \int_{E} f(x) \, dx.
\]
\end{namedres}
\begin{proof}
Sea $Z \subseteq E$ el conjunto nulo donde falla la convergencia, y $\tilde{E} = E \setminus Z$. Para cualquier función medible no negativa $g$ se tiene $R(g,Z) \subseteq Z \times [0,\infty) = \bigcup_{n} Z \times [0,n]$, que tiene medida cero; entonces $\int_{Z} g \, dx = 0$ y, por la aditividad del dominio, las integrales sobre $E$ y sobre $\tilde{E}$ coinciden tanto para $f$ como para cada $f_{k}$. Podemos entonces suponer que $f_{k}(x) \to f(x)$ para todo $x \in E$.

Con el mismo argumento de aproximación usado para la región bajo el gráfico, tenemos que
\[
R(f,E) = \Gamma(f,E) \cup \bigcup_{k=1}^{\infty} R(f_{k},E):
\]
si $0 \leq y < f(x)$, como $f_{k}(x) \uparrow f(x)$, existe $k$ con $y \leq f_{k}(x)$, es decir $(x,y) \in R(f_{k},E)$; si $y = f(x) < \infty$, entonces $(x,y) \in \Gamma(f,E)$; y cada $R(f_{k},E) \subseteq R(f,E)$ porque $f_{k} \leq f$. Como
\[
R(f_{k},E) \subseteq R(f_{k+1},E)
\]
por la monotonía de la sucesión, y $m(\Gamma(f,E)) = 0$, la continuidad desde abajo de la medida nos da
\[
\int_{E} f(x) \, dx = m(R(f,E)) = m\left( \bigcup_{k=1}^{\infty} R(f_{k},E) \right) = \lim_{k \to \infty} m(R(f_{k},E)) = \lim_{k \to \infty} \int_{E} f_{k}(x) \, dx. \qedhere
\]
\end{proof}

\subsection{Una fórmula útil}

\begin{namedstd}{Notación}{Suma inferior asociada a una partición finita del dominio}
Dada $f : E \to [0,\infty]$ medible y una partición finita $E = \bigcup_{i=1}^{N} E_{i}$ en medibles disjuntos dos a dos, definimos
\[
\sum\big(f, \{E_{i}\}_{i=1}^{N}\big) = \sum_{i=1}^{N} \Big( \inf_{E_{i}} f \Big) \ind_{E_{i}} \leq f.
\]
\end{namedstd}

\begin{namedres}{Teorema}{La integral como supremo de sumas inferiores}
Sea $f : E \to [0,\infty]$ medible, no negativa. Entonces
\[
\int_{E} f(x) \, dx = \sup \left\{ \int_{E} \sum\big(f, \{E_{i}\}_{i=1}^{N}\big) \, dx :\ E = \bigcup_{i=1}^{N} E_{i} \ \text{partición finita en medibles disjuntos} \right\}.
\]
\end{namedres}
\begin{proof}
Dados $E_{1}, \dots, E_{N}$ medibles disjuntos dos a dos tales que $E = \bigcup_{i=1}^{N} E_{i}$, tenemos que
\[
\sum\big(f, \{E_{i}\}_{i=1}^{N}\big) \leq f,
\]
entonces, por monotonía y la fórmula para funciones simples,
\[
\sum_{i=1}^{N} \Big( \inf_{E_{i}} f \Big) m(E_{i}) = \int_{E} \sum\big(f, \{E_{i}\}_{i=1}^{N}\big) \, dx \leq \int_{E} f(x) \, dx.
\]
Es decir, el supremo del enunciado es menor o igual que la integral.

Para la desigualdad recíproca, dados $k \geq 1$ y $1 \leq i \leq k 2^{k}$, defina
\[
E_{k}^{0} = \{ f \geq k \}, \qquad E_{k}^{i} = \left\{ \frac{i-1}{2^{k}} \leq f(x) < \frac{i}{2^{k}} \right\}.
\]
Entonces
\[
E = \bigcup_{i=0}^{k 2^{k}} E_{k}^{i}
\]
es una partición finita en medibles disjuntos dos a dos. Tome
\[
f_{k}(x) = k \ind_{E_{k}^{0}} + \sum_{i=1}^{k 2^{k}} \left( \frac{i-1}{2^{k}} \right) \ind_{E_{k}^{i}};
\]
esta es la aproximación diádica estándar, así que sabemos que
\[
f_{k} \leq f_{k+1} \qquad \text{y} \qquad \lim_{k \to \infty} f_{k} = f
\]
puntualmente en $E$. Por el teorema de convergencia monótona,
\[
\int_{E} f_{k} \, dx \to \int_{E} f \, dx.
\]
Como en cada pieza de la partición el valor de $f_{k}$ es menor o igual que el ínfimo de $f$ en ella, tenemos que
\[
f_{k} \leq \sum\big(f, \{E_{k}^{i}\}_{i=0}^{k 2^{k}}\big) \leq f,
\]
e integrando,
\[
\int_{E} f_{k} \, dx \leq \sum_{i=0}^{k 2^{k}} \Big( \inf_{E_{k}^{i}} f \Big) m(E_{k}^{i}) \leq \int_{E} f(x) \, dx.
\]
El término del medio es una de las sumas inferiores del enunciado, y el lado izquierdo converge a $\int_{E} f \, dx$; tenemos entonces que
\[
\lim_{k \to \infty} \sum_{i=0}^{k 2^{k}} \Big( \inf_{E_{k}^{i}} f \Big) m(E_{k}^{i}) = \int_{E} f(x) \, dx,
\]
y el supremo del enunciado alcanza a la integral.
\end{proof}

\begin{namedres}{Corolario}{La integral como supremo sobre funciones simples}
Sea $f : E \to [0,\infty]$ medible. Entonces
\[
\int_{E} f(x) \, dx = \sup \left\{ \int_{E} \phi(x) \, dx :\ \phi\ \text{simple y medible},\ 0 \leq \phi \leq f\ \text{en } E \right\}.
\]
\end{namedres}
\begin{proof}
Ejercicio.
\end{proof}

\subsection{Conjuntos nulos, desigualdad de Markov y funciones nulas}

\begin{namedres}{Lema}{La integral sobre un conjunto de medida cero es nula}
Sea $f : E \to [0,\infty]$ medible. Si $m(E) = 0$, entonces
\[
\int_{E} f(x) \, dx = 0.
\]
\end{namedres}
\begin{proof}
Sean $f_{k} = \sum_{i=1}^{m} a_{i} \ind_{A_{i}}$, con $A_{i} \cap A_{j} = \emptyset$ si $i \neq j$, funciones simples no negativas crecientes hacia $f$, de modo que, por el teorema de convergencia monótona,
\[
\lim_{k \to \infty} \int_{E} f_{k} \, dx = \int_{E} f(x) \, dx.
\]
Note que
\[
\int_{E} f_{k}(x) \, dx = \sum_{i=1}^{m} a_{i}\, m(A_{i}).
\]
Como $A_{i} \subseteq E$, entonces $m(A_{i}) = 0$ para $1 \leq i \leq m$, y cada $\int_{E} f_{k} \, dx = 0$. Concluimos que $\int_{E} f \, dx = 0$.
\end{proof}

\begin{namedres}{Teorema}{La integral respeta desigualdades c.p.d.}
Sean $f, g : E \to [0,\infty]$ medibles tales que $g(x) \leq f(x)$ casi por doquier en $E$. Entonces vale que
\[
\int_{E} g(x) \, dx \leq \int_{E} f(x) \, dx.
\]
En particular, si $f = g$ c.p.d. en $E$, entonces
\[
\int_{E} g(x) \, dx = \int_{E} f(x) \, dx.
\]
\end{namedres}
\begin{proof}
Sea
\[
A = \{ x \in E :\ g(x) \leq f(x) \}.
\]
Entonces $E = A \cup Z$ con $Z = E \setminus A$ de medida cero y $A \cap Z = \emptyset$. Así, por la aditividad del dominio y el lema anterior, tenemos
\[
\int_{E} f(x) \, dx = \int_{A} f(x) \, dx + \int_{Z} f(x) \, dx = \int_{A} f(x) \, dx,
\]
\[
\int_{E} g(x) \, dx = \int_{A} g(x) \, dx + \int_{Z} g(x) \, dx = \int_{A} g(x) \, dx.
\]
Como $g \leq f$ puntualmente en $A$, la monotonía de la integral da $\int_{A} g \, dx \leq \int_{A} f \, dx$, y el resultado se sigue. Si además $f = g$ c.p.d., ambas desigualdades valen y se obtiene la igualdad.
\end{proof}

\begin{namedres}{Proposición}{Desigualdad de Markov}
Sean $f : E \to [0,\infty]$ medible y $\alpha > 0$. Entonces
\[
m(\{ f \geq \alpha \}) \leq \frac{1}{\alpha} \int_{E} f(x) \, dx.
\]
\end{namedres}
\begin{proof}
Como $f \geq 0$, tenemos las desigualdades puntuales
\[
\alpha \ind_{\{f \geq \alpha\}} \leq f \ind_{\{f \geq \alpha\}} \leq f.
\]
Integrando y usando la fórmula para funciones simples, la identidad $\int_{E} \ind_{E_{1}} f \, dx = \int_{E_{1}} f \, dx$ y la monotonía,
\[
\alpha\, m(\{ f \geq \alpha \}) \leq \int_{\{f \geq \alpha\}} f(x) \, dx \leq \int_{E} f(x) \, dx,
\]
y basta dividir por $\alpha > 0$.
\end{proof}

\begin{namedres}{Corolario}{Una función no negativa con integral nula es nula c.p.d.}
Sea $f : E \to [0,\infty]$ medible tal que $\int_{E} f(x) \, dx = 0$. Entonces $f = 0$ casi por doquier.
\end{namedres}
\begin{proof}
Por la desigualdad de Markov, para todo $\alpha > 0$ tenemos que
\[
m(\{ f \geq \alpha \}) \leq \frac{1}{\alpha} \int_{E} f(x) \, dx = 0.
\]
Luego
\[
\{ f > 0 \} = \bigcup_{k=1}^{\infty} \left\{ f \geq \frac{1}{k} \right\}
\]
es una unión numerable de conjuntos de medida cero, y por lo tanto es un conjunto de medida cero. Es decir, $f = 0$ c.p.d.
\end{proof}

\subsection{Linealidad}

\begin{namedres}{Teorema}{Linealidad de la integral para funciones no negativas}
Sean $f, g : E \to [0,\infty]$ medibles y $c \geq 0$. Entonces
\[
\int_{E} \big( c f(x) + g(x) \big) \, dx = c \int_{E} f(x) \, dx + \int_{E} g(x) \, dx.
\]
\end{namedres}
\begin{proof}
\textit{Homogeneidad.} Si $c = 0$, con la convención $0 \cdot \infty = 0$ tenemos $cf \equiv 0$ y ambos lados de la homogeneidad son nulos. Sea entonces $c > 0$; vamos a probar que
\[
\int_{E} c f(x) \, dx = c \int_{E} f(x) \, dx.
\]
Sea $\{f_{k}\}$ una sucesión de funciones simples no negativas tales que
\[
0 \leq f_{k} \leq f_{k+1} \qquad \text{y} \qquad \lim_{k \to \infty} f_{k} = f.
\]
Si
\[
f_{k} = \sum_{i=1}^{m_{k}} a_{i} \ind_{A_{i}},
\]
con los $A_{i}$ medibles y disjuntos dos a dos, tenemos que
\[
\lim_{k \to \infty} c f_{k} = \lim_{k \to \infty} \sum_{i=1}^{m_{k}} c a_{i} \ind_{A_{i}} = c f(x), \qquad c f_{k} \leq c f_{k+1}.
\]
Entonces, aplicando el teorema de convergencia monótona dos veces y la fórmula para simples,
\[
\int_{E} c f(x) \, dx = \lim_{k \to \infty} \int_{E} c f_{k}(x) \, dx = \lim_{k \to \infty} \sum_{i=1}^{m_{k}} c a_{i}\, m(A_{i}) = c \lim_{k \to \infty} \sum_{i=1}^{m_{k}} a_{i}\, m(A_{i}) = c \int_{E} f(x) \, dx.
\]

\textit{Aditividad.} Falta mostrar la aditividad. Para tal efecto tomemos $g_{k}$ una sucesión de funciones simples no negativas tales que
\[
g_{k} \leq g_{k+1}, \qquad \lim_{k \to \infty} g_{k} = g.
\]
Entonces se sigue que
\[
g_{k} + f_{k} \leq g_{k+1} + f_{k+1}, \qquad \lim_{k \to \infty} (g_{k} + f_{k}) = g + f.
\]
Vamos a mostrar primero la aditividad para las simples, es decir, que
\[
\int_{E} \big( g_{k}(x) + f_{k}(x) \big) \, dx = \int_{E} g_{k}(x) \, dx + \int_{E} f_{k}(x) \, dx.
\]
Sea
\[
g_{k} = \sum_{j=1}^{\ell_{k}} b_{j} \ind_{B_{j}}, \quad \text{con } E = \bigcup_{j=1}^{\ell_{k}} B_{j} \text{ disjuntos dos a dos},
\]
y recordemos que
\[
f_{k} = \sum_{i=1}^{m_{k}} a_{i} \ind_{A_{i}}, \quad \text{con } E = \bigcup_{i=1}^{m_{k}} A_{i} \text{ disjuntos dos a dos}.
\]
Note que, al ser ambas familias particiones de $E$,
\[
\ind_{A_{i}} = \sum_{j=1}^{\ell_{k}} \ind_{B_{j} \cap A_{i}}, \qquad \ind_{B_{j}} = \sum_{i=1}^{m_{k}} \ind_{B_{j} \cap A_{i}},
\]
entonces
\[
g_{k}(x) + f_{k}(x) = \sum_{j=1}^{\ell_{k}} \sum_{i=1}^{m_{k}} b_{j} \ind_{A_{i} \cap B_{j}} + \sum_{i=1}^{m_{k}} \sum_{j=1}^{\ell_{k}} a_{i} \ind_{A_{i} \cap B_{j}} = \sum_{j=1}^{\ell_{k}} \sum_{i=1}^{m_{k}} (b_{j} + a_{i}) \ind_{A_{i} \cap B_{j}},
\]
que es una función simple sobre la partición común $\{A_{i} \cap B_{j}\}$. Ahora, por aditividad finita de la medida,
\[
\sum_{i=1}^{m_{k}} m(A_{i} \cap B_{j}) = m\left( B_{j} \cap \bigcup_{i=1}^{m_{k}} A_{i} \right) = m(B_{j} \cap E) = m(B_{j}),
\]
y de igual forma tenemos
\[
\sum_{j=1}^{\ell_{k}} m(A_{i} \cap B_{j}) = m(A_{i}).
\]
Finalmente,
\begin{align*}
\int_{E} \big( f_{k}(x) + g_{k}(x) \big) \, dx &= \sum_{j=1}^{\ell_{k}} \sum_{i=1}^{m_{k}} (a_{i} + b_{j})\, m(A_{i} \cap B_{j}) \\
&= \sum_{j=1}^{\ell_{k}} b_{j} \sum_{i=1}^{m_{k}} m(A_{i} \cap B_{j}) + \sum_{i=1}^{m_{k}} a_{i} \sum_{j=1}^{\ell_{k}} m(A_{i} \cap B_{j}) \\
&= \sum_{j=1}^{\ell_{k}} b_{j}\, m(B_{j}) + \sum_{i=1}^{m_{k}} a_{i}\, m(A_{i}) \\
&= \int_{E} g_{k}(x) \, dx + \int_{E} f_{k}(x) \, dx.
\end{align*}
Como $\{f_{k} + g_{k}\}$ es una sucesión creciente de simples que converge a $f + g$, el teorema de convergencia monótona aplicado tres veces nos da
\[
\int_{E} (f + g) \, dx = \lim_{k \to \infty} \int_{E} (f_{k} + g_{k}) \, dx = \lim_{k \to \infty} \left( \int_{E} f_{k} \, dx + \int_{E} g_{k} \, dx \right) = \int_{E} f \, dx + \int_{E} g \, dx.
\]
Combinando la homogeneidad con la aditividad, aplicadas a $cf$ y $g$, se concluye el enunciado.
\end{proof}

% =============================================================================
\section{La integral de Lebesgue de funciones medibles}
% =============================================================================

\subsection{La definición y las funciones integrables}

\begin{namedstd}{Definición}{Parte positiva, parte negativa y la integral general}
Sea $f : E \to \overline{\R}$ una función medible. Entonces
\[
f = f^{+} - f^{-}
\]
con
\[
f^{+} = \max(0, f) \qquad \text{y} \qquad f^{-} = \max(0, -f).
\]
Sabemos que $f^{+}$ y $f^{-}$ son medibles y no negativas, entonces
\[
\int_{E} f^{+} \, dx, \qquad \int_{E} f^{-} \, dx
\]
están bien definidas. Por tanto podemos definir
\[
\int_{E} f \, dx = \int_{E} f^{+} \, dx - \int_{E} f^{-} \, dx
\]
siempre que
\[
\int_{E} f^{+} \, dx < \infty \qquad \text{o} \qquad \int_{E} f^{-} \, dx < \infty.
\]
En tal caso decimos que la integral de $f$ sobre $E$ \emph{existe} (con valor posiblemente $\pm\infty$).
\end{namedstd}

\begin{namedstd}{Definición}{Función integrable y el espacio $L(E)$}
Decimos que $f$ es \emph{integrable}, o bien que $f \in L(E)$, si
\[
\int_{E} f^{+} \, dx < \infty \qquad \text{y} \qquad \int_{E} f^{-} \, dx < \infty.
\]
\end{namedstd}

\begin{namedstd}{Nota}{Desigualdad triangular integral y finitud c.p.d.}
Como $|f| = f^{+} + f^{-}$, para $f \in L(E)$ tenemos que
\[
\left| \int_{E} f \, dx \right| \leq \int_{E} f^{+} \, dx + \int_{E} f^{-} \, dx = \int_{E} |f| \, dx,
\]
donde la última igualdad usa la linealidad para funciones no negativas. En particular, $f \in L(E)$ si y solo si $\int_{E} |f| \, dx < \infty$. Además, si $\int_{E} |f| \, dx < \infty$, entonces $f \in \R$ c.p.d., por la proposición sobre funciones con integral finita.
\end{namedstd}

\subsection{Monotonía y propiedades básicas}

\begin{namedres}{Teorema}{Monotonía de la integral general}
Sean $f : E \to \overline{\R}$ y $g : E \to \overline{\R}$ medibles tales que
\[
\int_{E} f \, dx \qquad \text{y} \qquad \int_{E} g \, dx
\]
existen. Si $f \leq g$ c.p.d. en $E$, entonces
\[
\int_{E} f \, dx \leq \int_{E} g \, dx.
\]
En particular, si $f = g$ c.p.d. en $E$, entonces
\[
\int_{E} f \, dx = \int_{E} g \, dx.
\]
\end{namedres}
\begin{proof}
Note que si $f \leq g$ c.p.d., entonces, punto a punto donde vale la desigualdad,
\[
f^{+} = \max\{0, f\} \leq \max\{0, g\} = g^{+}, \qquad g^{-} = \max\{0, -g\} \leq \max\{0, -f\} = f^{-},
\]
es decir, $f^{+} \leq g^{+}$ y $g^{-} \leq f^{-}$ c.p.d. Luego, por la monotonía c.p.d. de la integral de funciones no negativas,
\[
\int_{E} f^{+} \, dx \leq \int_{E} g^{+} \, dx, \qquad \int_{E} g^{-} \, dx \leq \int_{E} f^{-} \, dx,
\]
y entonces, restando en el sentido extendido (la existencia de ambas integrales garantiza que no aparece la forma $\infty - \infty$),
\[
\int_{E} f \, dx = \int_{E} f^{+} \, dx - \int_{E} f^{-} \, dx \leq \int_{E} g^{+} \, dx - \int_{E} g^{-} \, dx = \int_{E} g \, dx.
\]
Si $f = g$ c.p.d., las dos desigualdades dan la igualdad.
\end{proof}

\begin{namedres}{Teorema}{Restricción del dominio, aditividad numerable y conjuntos nulos}
Sea $f : E \to \overline{\R}$ medible tal que $\int_{E} f \, dx$ existe.
\begin{enumerate}[label=(\roman*)]
\item Si $E_{1} \subseteq E$ es medible, entonces $\int_{E_{1}} f \, dx$ existe.
\item Si $E = \bigcup_{k=1}^{\infty} E_{k}$ con los $E_{k}$ medibles y disjuntos dos a dos, entonces
\[
\int_{E} f \, dx = \sum_{k=1}^{\infty} \int_{E_{k}} f \, dx.
\]
\item Si $E_{1} \subseteq E$ con $m(E_{1}) = 0$, entonces $\int_{E_{1}} f \, dx = 0$.
\end{enumerate}
\end{namedres}
\begin{proof}
Para (i), note que, por monotonía respecto al dominio para funciones no negativas,
\[
0 \leq \int_{E_{1}} f^{+} \, dx \leq \int_{E} f^{+} \, dx, \qquad 0 \leq \int_{E_{1}} f^{-} \, dx \leq \int_{E} f^{-} \, dx.
\]
Como alguna de las dos integrales sobre $E$ es finita, la correspondiente sobre $E_{1}$ también lo es, y $\int_{E_{1}} f \, dx$ existe.

Para (ii), sabemos, por la aditividad del dominio para funciones no negativas, que
\[
\int_{E} f^{+} \, dx = \sum_{k=1}^{\infty} \int_{E_{k}} f^{+} \, dx, \qquad \int_{E} f^{-} \, dx = \sum_{k=1}^{\infty} \int_{E_{k}} f^{-} \, dx.
\]
Como
\[
\int_{E} f^{+} \, dx = \sum_{k=1}^{\infty} \int_{E_{k}} f^{+} \, dx < \infty \qquad \text{o} \qquad \int_{E} f^{-} \, dx = \sum_{k=1}^{\infty} \int_{E_{k}} f^{-} \, dx < \infty,
\]
al menos una de las dos series converge, y podemos restar término a término sin ambigüedad. Entonces
\[
\int_{E} f \, dx = \sum_{k=1}^{\infty} \int_{E_{k}} f^{+} \, dx - \sum_{k=1}^{\infty} \int_{E_{k}} f^{-} \, dx = \sum_{k=1}^{\infty} \left( \int_{E_{k}} f^{+} \, dx - \int_{E_{k}} f^{-} \, dx \right) = \sum_{k=1}^{\infty} \int_{E_{k}} f \, dx.
\]

Finalmente, para (iii), si $m(E_{1}) = 0$, el lema de la integral sobre conjuntos nulos aplicado a $f^{+}$ y $f^{-}$ da
\[
\int_{E_{1}} f \, dx = \int_{E_{1}} f^{+} \, dx - \int_{E_{1}} f^{-} \, dx = 0 - 0 = 0. \qedhere
\]
\end{proof}

\subsection{Linealidad y diferencias de funciones}

\begin{namedres}{Teorema}{Linealidad de la integral general}
Sean $f, g : E \to \overline{\R}$ medibles tales que $\int_{E} f \, dx$ existe, y $c \in \R$. Entonces
\[
\int_{E} c f \, dx = c \int_{E} f \, dx.
\]
Además, si $f, g \in L(E)$, entonces $f + g \in L(E)$ y
\[
\int_{E} (f + g) \, dx = \int_{E} f \, dx + \int_{E} g \, dx.
\]
\end{namedres}
\begin{proof}
\textit{Homogeneidad.} Si $c \geq 0$, entonces $(cf)^{+} = c f^{+}$ y $(cf)^{-} = c f^{-}$, y el resultado se sigue de la homogeneidad para funciones no negativas. Si $c < 0$, entonces
\[
(cf)^{+} = -c f^{-}, \qquad (cf)^{-} = -c f^{+},
\]
pues $cf \geq 0$ exactamente donde $f \leq 0$. Entonces
\[
\int_{E} c f \, dx = \int_{E} (-c) f^{-} \, dx - \int_{E} (-c) f^{+} \, dx = -c \int_{E} f^{-} \, dx + c \int_{E} f^{+} \, dx = c \int_{E} f \, dx.
\]

\textit{Aditividad.} Si $f, g$ son integrables, entonces son finitas c.p.d., de modo que $f + g$ está definida c.p.d., y como $|f + g| \leq |f| + |g|$,
\[
0 \leq \int_{E} |f + g| \, dx \leq \int_{E} |f| \, dx + \int_{E} |g| \, dx < \infty,
\]
usando la monotonía y la linealidad para no negativas; es decir, $f + g \in L(E)$. Note que los siguientes conjuntos medibles, disjuntos dos a dos, cubren $E$ salvo un conjunto nulo (donde $f$ o $g$ toman valores infinitos):
\begin{align*}
&E_{1} = \{ f \geq 0,\ g \geq 0 \}, \\
&E_{2} = \{ f < 0,\ g < 0 \}, \\
&E_{3} = \{ f \geq 0,\ g < 0,\ f + g \geq 0 \}, \\
&E_{4} = \{ f < 0,\ g \geq 0,\ f + g \geq 0 \}, \\
&E_{5} = \{ f \geq 0,\ g < 0,\ f + g < 0 \}, \\
&E_{6} = \{ f < 0,\ g \geq 0,\ f + g < 0 \}.
\end{align*}
Por la aditividad del dominio, basta probar la aditividad de la integral en cada $E_{i}$.

En $E_{1}$ las tres funciones $f, g, f+g$ son no negativas, así que, por la linealidad para no negativas,
\[
\int_{E_{1}} (f + g) \, dx = \int_{E_{1}} (f^{+} + g^{+}) \, dx = \int_{E_{1}} f^{+} \, dx + \int_{E_{1}} g^{+} \, dx = \int_{E_{1}} f \, dx + \int_{E_{1}} g \, dx,
\]
y también, en $E_{2}$, donde las tres son negativas,
\begin{align*}
\int_{E_{2}} (f + g) \, dx &= \int_{E_{2}} -(f^{-} + g^{-}) \, dx = -\int_{E_{2}} (f^{-} + g^{-}) \, dx \\
&= -\int_{E_{2}} f^{-} \, dx - \int_{E_{2}} g^{-} \, dx = \int_{E_{2}} f \, dx + \int_{E_{2}} g \, dx.
\end{align*}
En $E_{3}$ tenemos $f = (f + g) + (-g)$ con $f + g \geq 0$ y $-g > 0$; de esta manera, por la linealidad para no negativas,
\[
\int_{E_{3}} f \, dx = \int_{E_{3}} \big( (f+g) - g \big) \, dx = \int_{E_{3}} (f + g) \, dx + \int_{E_{3}} (-g) \, dx = \int_{E_{3}} (f + g) \, dx - \int_{E_{3}} g \, dx,
\]
donde el último paso usa la homogeneidad con $c = -1$ y que $g \in L(E_{3})$. Por lo tanto
\[
\int_{E_{3}} (f + g) \, dx = \int_{E_{3}} f \, dx + \int_{E_{3}} g \, dx.
\]
Además, en $E_{5}$ tenemos $-g = -(f+g) + f$ con $-(f+g) > 0$ y $f \geq 0$, entonces
\[
\int_{E_{5}} (-g) \, dx = \int_{E_{5}} \big( -(f+g) + f \big) \, dx = \int_{E_{5}} -(f+g) \, dx + \int_{E_{5}} f \, dx.
\]
Así, multiplicando por $-1$ y reordenando, tenemos que
\[
\int_{E_{5}} f \, dx + \int_{E_{5}} g \, dx = \int_{E_{5}} (f + g) \, dx.
\]
Los casos $E_{4}$ y $E_{6}$ son análogos, intercambiando los papeles de $f$ y $g$; terminar el resto de la prueba es un ejercicio.
\end{proof}

\begin{namedstd}{Nota}{Combinaciones lineales y una advertencia sobre la resta}
Sean $f_{1}, \dots, f_{n} \in L(E)$ y $a_{1}, \dots, a_{n} \in \R$. Entonces, iterando el teorema anterior, $\sum_{k=1}^{n} a_{k} f_{k} \in L(E)$ y
\[
\int_{E} \left( \sum_{k=1}^{n} a_{k} f_{k} \right) \, dx = \sum_{k=1}^{n} a_{k} \int_{E} f_{k} \, dx.
\]
Ojo: cuando las integrales no son finitas, en general \emph{no} es cierto que
\[
\int_{E} (f - g) \, dx = \int_{E} f \, dx - \int_{E} g \, dx.
\]
Por ejemplo, considere $f = \ind_{[n,\infty)}$ y $g = \ind_{[n+1,\infty)}$: la resta $f - g = \ind_{[n,n+1)}$ tiene integral $1$, pero $\int f \, dx = \int g \, dx = \infty$ y la diferencia $\infty - \infty$ no está definida.
\end{namedstd}

\begin{namedres}{Proposición}{Resta de una minorante integrable}
Sean $f$ y $\phi$ medibles en $E$ tales que $\phi \leq f$ c.p.d. y $\phi \in L(E)$. Entonces $\int_{E} f \, dx$ y $\int_{E} (f - \phi) \, dx$ existen y
\[
\int_{E} (f - \phi) \, dx = \int_{E} f \, dx - \int_{E} \phi \, dx.
\]
\end{namedres}
\begin{proof}
Como $\phi \leq f$ c.p.d., tenemos $-f \leq -\phi$ c.p.d. y entonces
\[
f^{-} = \max(0, -f) \leq \max(0, -\phi) = \phi^{-} \quad \text{c.p.d.},
\]
de donde
\[
0 \leq \int_{E} f^{-} \, dx \leq \int_{E} \phi^{-} \, dx < \infty.
\]
Es decir, $\int_{E} f \, dx$ existe, con valor en $(-\infty, \infty]$, y quedan dos casos según el valor de $\int_{E} f^{+} \, dx$.

Si $\int_{E} f^{+} \, dx < \infty$, entonces $f \in L(E)$, y como también $-\phi \in L(E)$, la linealidad da directamente
\[
\int_{E} (f - \phi) \, dx = \int_{E} f \, dx - \int_{E} \phi \, dx.
\]

Ahora, si $\int_{E} f^{+} \, dx = \infty$, es decir $f \notin L(E)$, afirmamos que $f - \phi \notin L(E)$. En efecto, donde ambas son finitas vale $f = \phi + (f - \phi) \leq \phi^{+} + (f - \phi)$, y como el lado derecho es no negativo,
\[
f^{+} \leq \phi^{+} + (f - \phi) \quad \text{c.p.d.}
\]
Integrando, por monotonía c.p.d. y linealidad para no negativas,
\[
\infty = \int_{E} f^{+} \, dx \leq \int_{E} \phi^{+} \, dx + \int_{E} (f - \phi) \, dx,
\]
y como $\int_{E} \phi^{+} \, dx < \infty$, se sigue que $\int_{E} (f - \phi) \, dx = \infty$. Además $f - \phi \geq 0$ c.p.d., así que
\[
\int_{E} (f - \phi) \, dx = \int_{E} (f - \phi)^{+} \, dx = \infty.
\]
Concluimos que, en este caso, ambos lados valen $\infty$:
\[
\int_{E} (f - \phi) \, dx = \infty = \int_{E} f \, dx - \int_{E} \phi \, dx,
\]
pues $\int_{E} f \, dx = \infty$ y $\int_{E} \phi \, dx \in \R$.
\end{proof}

\begin{namedstd}{Nota}{Productos de funciones integrables}
Las condiciones para que $fg$ sea integrable son más complejas que para la suma. Por ejemplo, si $f \in L(E)$ y $g$ es medible con $|g(x)| \leq M$ para $x \in E$, entonces $fg \in L(E)$, pues $|fg| \leq M |f|$ y la monotonía da $\int_{E} |fg| \, dx \leq M \int_{E} |f| \, dx < \infty$.
\end{namedstd}

\subsection{Los teoremas de convergencia}

\begin{namedres}{Teorema}{Convergencia monótona con minorante o mayorante integrable}
Sea $\{f_{k}\}_{k=1}^{\infty}$ una sucesión de funciones medibles en $E$, tales que $\lim_{k \to \infty} f_{k} = f$ c.p.d. en $E$.
\begin{enumerate}[label=(\roman*)]
\item Si existe $\phi \in L(E)$ tal que $\phi \leq f_{k} \leq f_{k+1}$ c.p.d. para todo $k \geq 1$, entonces
\[
\lim_{k \to \infty} \int_{E} f_{k} \, dx = \int_{E} f \, dx.
\]
\item Si existe $\phi \in L(E)$ tal que $f_{k+1} \leq f_{k} \leq \phi$ c.p.d. para todo $k \geq 1$, entonces
\[
\lim_{k \to \infty} \int_{E} f_{k} \, dx = \int_{E} f \, dx.
\]
\end{enumerate}
\end{namedres}
\begin{proof}
Solo probaremos el caso creciente; la parte (ii) se deja como ejercicio a la persona lectora. Si se cumple la primera condición, entonces
\[
f_{k} - \phi \geq 0 \qquad \text{y} \qquad f_{k+1} - \phi \geq f_{k} - \phi \quad \text{c.p.d.}
\]
para $k \geq 1$, y además $f_{k} - \phi \to f - \phi$ c.p.d. Entonces, por el teorema de convergencia monótona para funciones no negativas,
\[
\lim_{k \to \infty} \int_{E} (f_{k} - \phi) \, dx = \int_{E} (f - \phi) \, dx.
\]
Por la proposición de la resta de una minorante integrable, aplicada a cada $f_{k}$ y a $f$ (con minorante $\phi$),
\[
\int_{E} (f_{k} - \phi) \, dx = \int_{E} f_{k} \, dx - \int_{E} \phi \, dx, \qquad \int_{E} (f - \phi) \, dx = \int_{E} f \, dx - \int_{E} \phi \, dx.
\]
Es decir,
\[
\lim_{k \to \infty} \int_{E} f_{k} \, dx - \int_{E} \phi \, dx = \int_{E} f \, dx - \int_{E} \phi \, dx,
\]
y como $\int_{E} \phi \, dx \in \R$, podemos cancelar y concluir que $\lim_{k \to \infty} \int_{E} f_{k} \, dx = \int_{E} f \, dx$.
\end{proof}

\begin{namedstd}{Ejemplo}{Integración término a término de series de funciones no negativas}
Sean $f_{k} : E \to [0,\infty]$ medibles para $k \geq 1$. Entonces, si llamamos
\[
g_{m} = \sum_{k=1}^{m} f_{k},
\]
tenemos que $0 \leq g_{m} \leq g_{m+1}$. Luego, por el teorema de convergencia monótona y la linealidad,
\begin{align*}
\int_{E} \left( \sum_{k=1}^{\infty} f_{k} \right) \, dx &= \int_{E} \lim_{m \to \infty} g_{m}(x) \, dx = \lim_{m \to \infty} \int_{E} g_{m}(x) \, dx \\
&= \lim_{m \to \infty} \int_{E} \sum_{k=1}^{m} f_{k}(x) \, dx = \lim_{m \to \infty} \sum_{k=1}^{m} \int_{E} f_{k}(x) \, dx = \sum_{k=1}^{\infty} \int_{E} f_{k}(x) \, dx.
\end{align*}
\end{namedstd}

\begin{namedstd}{Ejemplo}{Una sucesión que converge a cero con integrales constantes}
Sea $E = [0,1]$ y definamos
\[
f_{k}(x) =
\begin{cases}
k & \text{si } 0 \leq x \leq \frac{1}{k}, \\
0 & \text{si } \frac{1}{k} < x \leq 1.
\end{cases}
\]
Entonces
\[
\int_{E} f_{k}(x) \, dx = k \cdot m\left( \left[0, \tfrac{1}{k}\right] \right) = 1 \quad \text{para todo } k,
\]
pero $\lim_{k \to \infty} f_{k} = 0$ c.p.d. en $[0,1]$: para $x > 0$, $f_{k}(x) = 0$ apenas $k > \tfrac{1}{x}$. Nos preguntamos:
\begin{center}
¿Cuándo $\displaystyle \lim_{k \to \infty} \int_{E} f_{k}(x) \, dx = \int_{E} \lim_{k \to \infty} f_{k}(x) \, dx$?
\end{center}
Este ejemplo muestra que alguna hipótesis adicional (monotonía, dominación, etc.) es necesaria.
\end{namedstd}

\begin{namedres}{Teorema}{Lema de Fatou}
Sea $f_{k} : E \to [0,\infty]$ una sucesión de funciones medibles no negativas. Entonces
\[
\int_{E} \liminf_{k \to \infty} f_{k} \, dx \leq \liminf_{k \to \infty} \int_{E} f_{k} \, dx.
\]
\end{namedres}
\begin{proof}
Recordemos que
\[
\liminf_{k \to \infty} f_{k} = \sup_{k \geq 1} \inf_{m \geq k} f_{m} = \lim_{k \to \infty} \inf_{m \geq k} f_{m}.
\]
Considere
\[
g_{k} = \inf_{m \geq k} f_{m};
\]
entonces $\{g_{k}\}$ es una sucesión creciente de funciones medibles no negativas con $g_{k} \to \liminf_{k \to \infty} f_{k}$. Por el teorema de convergencia monótona,
\[
\int_{E} \Big( \liminf_{k \to \infty} f_{k} \Big) \, dx = \lim_{k \to \infty} \int_{E} \inf_{m \geq k} f_{m} \, dx.
\]
Note que
\[
\inf_{m \geq k} f_{m} \leq f_{n} \quad \text{para } n \geq k \geq 1.
\]
Entonces, por monotonía, para $n \geq k$,
\[
\int_{E} \inf_{m \geq k} f_{m}(x) \, dx \leq \int_{E} f_{n}(x) \, dx,
\]
y tomando el ínfimo sobre $n \geq k$,
\[
\int_{E} \inf_{m \geq k} f_{m}(x) \, dx \leq \inf_{n \geq k} \int_{E} f_{n}(x) \, dx.
\]
Tomando ahora el límite cuando $k \to \infty$ en ambos lados,
\[
\lim_{k \to \infty} \int_{E} \inf_{m \geq k} f_{m}(x) \, dx \leq \lim_{k \to \infty} \inf_{n \geq k} \int_{E} f_{n}(x) \, dx = \liminf_{k \to \infty} \int_{E} f_{k}(x) \, dx.
\]
Combinando con la primera identidad, concluimos que
\[
\int_{E} \liminf_{k \to \infty} f_{k} \, dx \leq \liminf_{k \to \infty} \int_{E} f_{k} \, dx. \qedhere
\]
\end{proof}

\begin{namedstd}{Nota}{Cota uniforme de integrales y el límite inferior}
Note que si $f_{k} : E \to [0,\infty]$ es medible para cada $k$ y
\[
\int_{E} f_{k}(x) \, dx \leq M \quad \text{para todo } k,
\]
entonces, por el lema de Fatou,
\[
\int_{E} \liminf_{k \to \infty} f_{k} \, dx \leq M.
\]
\end{namedstd}

\begin{namedres}{Teorema}{Convergencia dominada de Lebesgue para funciones no negativas}
Sea $\{f_{k}\}_{k=1}^{\infty}$ una sucesión de funciones medibles no negativas tales que
\[
\lim_{k \to \infty} f_{k} = f
\]
c.p.d. en $E$. Si existe $\phi$ medible con
\[
0 \leq f_{k} \leq \phi \quad \text{c.p.d. para todo } k \geq 1, \qquad \int_{E} \phi(x) \, dx < \infty,
\]
entonces
\[
\lim_{k \to \infty} \int_{E} f_{k}(x) \, dx = \int_{E} f(x) \, dx.
\]
\end{namedres}
\begin{proof}
Como $f = \liminf_{k \to \infty} f_{k}$ c.p.d., el lema de Fatou y la igualdad c.p.d. de las integrales dan
\[
\int_{E} f(x) \, dx \leq \liminf_{k \to \infty} \int_{E} f_{k} \, dx.
\]
Note además que $0 \leq f \leq \phi$ c.p.d., entonces $f, f_{k} \in L(E)$, pues sus integrales están dominadas por $\int_{E} \phi \, dx < \infty$.

Considere ahora $h_{k} = \phi - f_{k} \geq 0$ c.p.d. Por el lema de Fatou aplicado a $\{h_{k}\}$,
\[
\int_{E} \liminf_{k \to \infty} (\phi - f_{k}) \, dx \leq \liminf_{k \to \infty} \int_{E} (\phi - f_{k}) \, dx.
\]
Note que, como el límite de $f_{k}$ existe c.p.d.,
\[
\liminf_{k \to \infty} (\phi - f_{k}) = \lim_{k \to \infty} (\phi - f_{k}) = \phi - f \quad \text{c.p.d.}
\]
Por otro lado, usando la linealidad (todas las funciones involucradas son integrables), tenemos que
\[
\liminf_{k \to \infty} \int_{E} (\phi - f_{k})(x) \, dx = \liminf_{k \to \infty} \left( \int_{E} \phi(x) \, dx - \int_{E} f_{k}(x) \, dx \right) = \int_{E} \phi(x) \, dx - \limsup_{k \to \infty} \int_{E} f_{k}(x) \, dx.
\]
La última igualdad vale porque $\liminf_{k}(-a_{k}) = -\limsup_{k} a_{k}$ para toda sucesión real $\{a_{k}\}$, y sumar la constante finita $\int_{E} \phi \, dx$ conmuta con el límite inferior. Combinando las tres relaciones anteriores y cancelando $\int_{E} \phi \, dx \in \R$, concluimos que
\[
\limsup_{k \to \infty} \int_{E} f_{k}(x) \, dx \leq \int_{E} f(x) \, dx \leq \liminf_{k \to \infty} \int_{E} f_{k}(x) \, dx,
\]
es decir, el límite existe y $\lim_{k \to \infty} \int_{E} f_{k} \, dx = \int_{E} f \, dx$.
\end{proof}

\begin{namedres}{Teorema}{Convergencia uniforme en dominios de medida finita}
Sea $f_{k} \in L(E)$ para $k \geq 1$. Si $\lim_{k \to \infty} f_{k} = f$ uniformemente en $E$ con $m(E) < \infty$, entonces $f \in L(E)$ y
\[
\lim_{k \to \infty} \int_{E} f_{k} \, dx = \int_{E} f \, dx.
\]
\end{namedres}
\begin{proof}
Ejercicio.
\end{proof}

\begin{namedstd}{Nota}{La hipótesis de medida finita es esencial en la convergencia uniforme}
Note que $f_{k}(x) = \frac{1}{k}$ converge a cero uniformemente en $\R$, pero
\[
\int_{\R} f_{k} \, dx = \infty
\]
para todo $k$, así que la conclusión del teorema anterior falla sin la hipótesis $m(E) < \infty$.
\end{namedstd}

% =============================================================================
\section{La relación entre las integrales de Riemann y Lebesgue}
% =============================================================================

\subsection{Particiones y funciones escalonadas}

\begin{namedstd}{Notación}{Particiones y funciones escalonadas asociadas}
Sea $f : [a,b] \to \R$ acotada por $M$, es decir $|f(x)| \leq M$ en $[a,b]$, y sea
\[
\Gamma_{k} = \{ a = x_{0}^{k} < x_{1}^{k} < \dots < x_{n_{k}}^{k} = b \}
\]
una sucesión de particiones de $[a,b]$ tales que $\Gamma_{k} \subseteq \Gamma_{k+1}$ (cada una refina a la anterior) y cuya malla $|\Gamma_{k}| = \max_{i} (x_{i}^{k} - x_{i-1}^{k})$ tiende a cero. Escribiendo $x_{i} = x_{i}^{k}$ y $n = n_{k}$ para aligerar la notación, consideremos:
\begin{enumerate}[label=(\alph*)]
\item si $m_{i} = \inf_{[x_{i-1},x_{i}]} f(x)$, la \emph{escalonada inferior}
\[
\ell_{k}(x) = \sum_{i=1}^{n-1} m_{i} \ind_{[x_{i-1},x_{i})} + m_{n} \ind_{[x_{n-1},x_{n}]};
\]
\item si $M_{i} = \sup_{[x_{i-1},x_{i}]} f(x)$, la \emph{escalonada superior}
\[
u_{k}(x) = \sum_{i=1}^{n-1} M_{i} \ind_{[x_{i-1},x_{i})} + M_{n} \ind_{[x_{n-1},x_{n}]}.
\]
\end{enumerate}
Las integrales de estas funciones simples son las sumas de Darboux:
\[
\int_{[a,b]} \ell_{k}(x) \, dx = \sum_{i=1}^{n} m_{i} (x_{i} - x_{i-1}), \qquad \int_{[a,b]} u_{k}(x) \, dx = \sum_{i=1}^{n} M_{i} (x_{i} - x_{i-1}).
\]
\end{namedstd}

\begin{namedstd}{Nota}{Monotonía de las escalonadas bajo refinamiento}
Note que
\[
\ell_{k}(x) \leq f(x) \leq u_{k}(x) \quad \text{para todo } x \in [a,b].
\]
Recordemos que si $\Gamma_{k} \subseteq \Gamma_{k+1}$, entonces
\begin{enumerate}[label=(\roman*)]
\item $\ell_{k} \leq \ell_{k+1}$;
\item $u_{k} \geq u_{k+1}$.
\end{enumerate}
Como $|f(x)| \leq M$ para $x \in [a,b]$, entonces $|\ell_{k}| \leq M$ y $|u_{k}| \leq M$ en $[a,b]$. Por monotonía y acotación, existen los límites puntuales
\[
\ell = \lim_{k \to \infty} \ell_{k}, \qquad u = \lim_{k \to \infty} u_{k},
\]
que son medibles por ser límites de medibles, y satisfacen $\ell(x) \leq f(x) \leq u(x)$ para todo $x \in [a,b]$.
\end{namedstd}

\begin{namedres}{Proposición}{Las sumas de Darboux convergen a las integrales de $\ell$ y $u$}
Con la notación anterior,
\[
\lim_{k \to \infty} \int_{[a,b]} \ell_{k}(x) \, dx = \int_{[a,b]} \ell(x) \, dx, \qquad \lim_{k \to \infty} \int_{[a,b]} u_{k}(x) \, dx = \int_{[a,b]} u(x) \, dx.
\]
Es decir, estas integrales son el límite de las sumas inferiores y superiores de Darboux. Además,
\[
\int_{[a,b]} u(x) \, dx = \int_{[a,b]} \ell(x) \, dx \iff \int_{[a,b]} \big( u(x) - \ell(x) \big) \, dx = 0 \iff u = \ell \ \text{c.p.d.}
\]
\end{namedres}
\begin{proof}
Las funciones $\ell_{k} + M$ y $u_{k} + M$ son medibles, no negativas, convergen puntualmente a $\ell + M$ y $u + M$ respectivamente, y están dominadas por la constante $2M$, que es integrable en $[a,b]$ pues $m([a,b]) = b - a < \infty$. El teorema de convergencia dominada implica entonces que
\[
\lim_{k \to \infty} \int_{[a,b]} (\ell_{k} + M) \, dx = \int_{[a,b]} (\ell + M) \, dx,
\]
y restando la constante $M(b-a)$ de ambos lados (linealidad) se obtiene la afirmación para $\ell_{k}$; el argumento para $u_{k}$ es idéntico.

Para las equivalencias: como $u - \ell \geq 0$ y ambas son integrables (están acotadas por $M$ en un dominio de medida finita), la linealidad da
\[
\int_{[a,b]} u \, dx - \int_{[a,b]} \ell \, dx = \int_{[a,b]} (u - \ell) \, dx,
\]
así que la primera equivalencia es inmediata. La segunda es el corolario de la desigualdad de Markov: una función no negativa tiene integral nula si y solo si es nula c.p.d.
\end{proof}

\begin{namedstd}{Nota}{Si las sumas coinciden, la integral de Riemann es la de Lebesgue}
Concluimos que, si las sumas inferiores y superiores de Darboux convergen al mismo valor, entonces
\[
u = \ell = f \quad \text{c.p.d.}
\]
(pues $\ell \leq f \leq u$), la función $f$ es medible por coincidir c.p.d. con la función medible $\ell$, y por lo tanto
\[
\lim_{k \to \infty} \sum_{i=1}^{n} m_{i} (x_{i} - x_{i-1}) = \lim_{k \to \infty} \sum_{i=1}^{n} M_{i} (x_{i} - x_{i-1}) = \int_{[a,b]} f(x) \, dx,
\]
donde la integral es la de Lebesgue. Es decir, cuando $f$ es Riemann integrable, su integral de Riemann coincide con su integral de Lebesgue.
\end{namedstd}

\subsection{La caracterización de Lebesgue de la integrabilidad de Riemann}

\begin{namedres}{Teorema}{Criterio de Lebesgue: Riemann integrable equivale a continua c.p.d.}
Sea $f : [a,b] \to \R$ acotada. Entonces son equivalentes:
\begin{enumerate}[label=(\roman*)]
\item $f$ es Riemann integrable;
\item $f$ es continua c.p.d. en $[a,b]$.
\end{enumerate}
En tal caso, las integrales de Riemann y de Lebesgue de $f$ sobre $[a,b]$ coinciden.
\end{namedres}
\begin{proof}
Fijamos una sucesión de particiones $\Gamma_{k}$ con $\Gamma_{k} \subseteq \Gamma_{k+1}$ y $\lim_{k \to \infty} |\Gamma_{k}| = 0$, y usamos la notación de la subsección anterior.

\textit{(i) implica (ii).} Comenzamos asumiendo que $f$ es Riemann integrable; entonces las sumas inferiores y superiores convergen ambas a la integral de Riemann, y por la proposición anterior $u = \ell = f$ c.p.d. Sea
\[
Z = \{ \ell \neq f \} \cup \{ \ell \neq u \} \cup \{ f \neq u \} \cup \bigcup_{k=1}^{\infty} \Gamma_{k}.
\]
Entonces $m(Z) = 0$, pues es unión numerable de conjuntos nulos: los tres primeros por lo anterior, y cada $\Gamma_{k}$ por ser finito. Si $x \notin Z$, entonces
\[
u(x) = \ell(x) = f(x), \qquad x \notin \Gamma_{k} \ \text{para todo } k \geq 1.
\]
Afirmamos que $f$ es continua en cada $x \notin Z$. Supongamos que $f$ no es continua en $x$. Entonces existe un $\varepsilon > 0$ tal que para todo $\delta > 0$ existe $x_{\delta}$ con
\[
|x_{\delta} - x| < \delta \qquad \text{y} \qquad |f(x) - f(x_{\delta})| > \varepsilon.
\]
Dado $k$, como $x \notin \Gamma_{k}$, existen nodos consecutivos $x_{i-1}^{k}, x_{i}^{k}$ tales que $x \in (x_{i-1}^{k}, x_{i}^{k})$, que es abierto. Tomemos $\delta > 0$ tal que
\[
(x - \delta, x + \delta) \subseteq (x_{i-1}^{k}, x_{i}^{k}).
\]
Entonces $x_{\delta} \in (x_{i-1}^{k}, x_{i}^{k})$ y, como ambos puntos pertenecen al mismo subintervalo de la partición,
\[
\varepsilon < |f(x) - f(x_{\delta})| \leq M_{i} - m_{i} = u_{k}(x) - \ell_{k}(x).
\]
Como esto vale para todo $k$, haciendo $k \to \infty$ obtenemos
\[
u(x) - \ell(x) \geq \varepsilon > 0,
\]
lo que contradice $u(x) = \ell(x)$. Concluimos que $f$ es continua en todo punto fuera del conjunto nulo $Z$, es decir, $f$ es continua c.p.d.

\textit{(ii) implica (i).} Asumamos ahora que
\[
Z = \{ x \in [a,b] :\ f\ \text{es discontinua en } x \}
\]
es de medida cero. Sean $x \notin Z \cup \{a, b\}$ y $\varepsilon > 0$. Por continuidad de $f$ en $x$, existe $\delta > 0$ tal que
\[
|x - y| < \delta \implies |f(x) - f(y)| < \frac{\varepsilon}{2}.
\]
Tome ahora $k_{0}$ tal que
\[
|\Gamma_{k}| < \frac{\delta}{2} \quad \text{si } k \geq k_{0}.
\]
Para $k \geq k_{0}$, sea $i$ tal que $x \in [x_{i-1}, x_{i})$. Entonces todo punto $y$ del subintervalo cerrado $[x_{i-1}, x_{i}]$ satisface $|y - x| \leq x_{i} - x_{i-1} \leq |\Gamma_{k}| < \delta$, es decir,
\[
[x_{i-1}, x_{i}] \subseteq (x - \delta, x + \delta).
\]
Por lo tanto,
\[
f(x) - \frac{\varepsilon}{2} \leq f(y) \leq f(x) + \frac{\varepsilon}{2} \quad \text{para } y \in [x_{i-1}, x_{i}],
\]
y tomando ínfimo y supremo sobre ese subintervalo,
\[
f(x) - \frac{\varepsilon}{2} \leq m_{i} \leq M_{i} \leq f(x) + \frac{\varepsilon}{2}.
\]
Es decir, para $k \geq k_{0}$,
\[
|u_{k}(x) - f(x)| \leq \frac{\varepsilon}{2} < \varepsilon, \qquad |\ell_{k}(x) - f(x)| \leq \frac{\varepsilon}{2} < \varepsilon.
\]
Lo anterior nos dice que
\[
\lim_{k \to \infty} \ell_{k}(x) = f(x), \qquad \lim_{k \to \infty} u_{k}(x) = f(x)
\]
para todo $x$ fuera del conjunto nulo $Z \cup \{a,b\}$, es decir, $\ell = u = f$ c.p.d.; en particular $f$ es medible. Por el teorema de convergencia dominada (aplicado como en la proposición anterior, con dominación por la constante $M$ en el dominio de medida finita $[a,b]$), se cumple que
\[
\lim_{k \to \infty} \int_{[a,b]} \ell_{k}(x) \, dx = \int_{[a,b]} f(x) \, dx, \qquad \lim_{k \to \infty} \int_{[a,b]} u_{k}(x) \, dx = \int_{[a,b]} f(x) \, dx.
\]
Es decir, las sumas inferiores y superiores de Darboux convergen ambas al mismo valor; en particular, dado $\varepsilon > 0$ existe una partición cuya suma superior y suma inferior difieren en menos de $\varepsilon$, y por el criterio de Darboux $f$ es Riemann integrable. En conclusión, $f$ es Riemann integrable y las integrales de Riemann y Lebesgue concuerdan.
\end{proof}

\appendix
% =============================================================================
\section{Anexo: caracterizaciones y teoremas clave}
% =============================================================================

A lo largo de este parcial, varios conceptos clave admiten \emph{múltiples caracterizaciones equivalentes} o vienen empaquetados en teoremas con hipótesis precisas que conviene tener juntas a la vista. Este anexo recopila esos paquetes en cajas autocontenidas, una por concepto. Cada caja indica además dónde se prueba cada afirmación dentro del documento.

\anexogroup{Egorov y Lusin}

\begin{equivbox}{blue}{Teorema de Egorov --- convergencia casi uniforme}
Sean $E \subseteq \R^{d}$ medible con $m(E) < \infty$, $\{f_{k}\}$ medibles, $f_{k} \to f$ c.p.d. y $|f| < \infty$ c.p.d. en $E$. Entonces, para todo $\varepsilon > 0$ existe $F_{\varepsilon} \subseteq E$ \emph{cerrado} con:
\begin{enumerate}[label=(\arabic*),leftmargin=2em]
\item $m(E \setminus F_{\varepsilon}) < \varepsilon$; \hfill\textit{Teorema 1.1.5.}
\item $f_{k} \to f$ uniformemente en $F_{\varepsilon}$.
\end{enumerate}
\textbf{Hipótesis esenciales:} $m(E) < \infty$ (contraejemplo: $f_{k} = \ind_{B(0,k)}$ en $\R^{d}$, Ejemplo 1.1.2) y $|f| < \infty$ c.p.d.
\end{equivbox}

\begin{equivbox}{teal}{Teorema de Lusin --- caracterización de la medibilidad}
Sea $f : E \to \R$ con $E$ medible. Las siguientes son equivalentes (Teorema 1.2.4):
\begin{enumerate}[label=(\arabic*),leftmargin=2em]
\item \textbf{(medibilidad)}\; $f$ es medible.
\item \textbf{(propiedad $C$)}\; Para todo $\varepsilon > 0$ existe $F \subseteq E$ cerrado con $m(E \setminus F) < \varepsilon$ y $f|_{F}$ continua. \hfill\textit{Definición 1.2.1.}
\end{enumerate}
\textbf{Cadena de la prueba:} simples tienen propiedad $C$ (Lema 1.2.3) $\to$ caso $m(E) < \infty$ vía Egorov $\to$ caso $m(E) = \infty$ vía anillos $E_{k} = E \cap \{k-1 \leq |x| < k\}$.
\end{equivbox}

\anexogroup{Modos de convergencia}

\begin{equivbox}{violet}{Relaciones entre los modos de convergencia}
Sean $f, f_{n} : E \to \overline{\R}$ medibles y finitas c.p.d.
\begin{itemize}[leftmargin=2em,itemsep=2pt]
\item Uniforme $\Rightarrow$ puntual $\Rightarrow$ c.p.d. (siempre).
\item C.p.d. $+$ $m(E) < \infty$ $\Rightarrow$ en medida. \hfill\textit{Teorema 2.1.3.}
\item En medida $\Rightarrow$ existe subsucesión que converge c.p.d. \hfill\textit{Teorema 2.1.5.}
\item Cauchy en medida $\Rightarrow$ convergente en medida. \hfill\textit{Lema 2.2.2 (ejercicio).}
\item C.p.d. $\not\Rightarrow$ en medida si $m(E) = \infty$ (e.g., $f_{k} = \ind_{[k,k+1]}$ en $\R$).
\item En medida $\not\Rightarrow$ c.p.d. (máquina de escribir, Ejemplo 2.1.4).
\item C.p.d. $\not\Rightarrow$ uniforme, pero Egorov da uniformidad fuera de un conjunto pequeño si $m(E) < \infty$.
\end{itemize}
\end{equivbox}

\anexogroup{La integral de funciones no negativas}

\begin{equivbox}{purple}{Tres caracterizaciones de $\int_{E} f \, dx$ para $f \geq 0$}
Sea $f : E \to [0,\infty]$ medible. Las siguientes cantidades coinciden:
\begin{enumerate}[label=(\arabic*),leftmargin=2em]
\item \textbf{(geométrica)}\; $m(R(f,E))$, la medida de la región bajo el gráfico. \hfill\textit{Definición 3.3.1.}
\item \textbf{(sumas inferiores)}\; $\sup \big\{ \sum_{i} (\inf_{E_{i}} f)\, m(E_{i}) :\ E = \bigcup_{i=1}^{N} E_{i}$ partición medible finita$\big\}$. \hfill\textit{Teorema 4.3.2.}
\item \textbf{(funciones simples)}\; $\sup \big\{ \int_{E} \phi \, dx :\ \phi$ simple medible, $0 \leq \phi \leq f \big\}$. \hfill\textit{Corolario 4.3.3.}
\end{enumerate}
\textbf{Insumos geométricos:} $m(A \times [0,a]) = a\,m(A)$ (Lema 3.1.3); $m_{e}(\Gamma(f,E)) = 0$ (Lema 3.2.1); $R(f,E)$ medible (Teorema 3.2.4).
\end{equivbox}

\begin{equivbox}{orange!85!black}{Propiedades básicas de la integral no negativa}
Sean $f, g : E \to [0,\infty]$ medibles y $c \geq 0$.
\begin{itemize}[leftmargin=2em,itemsep=2pt]
\item \textbf{Monotonía:} $g \leq f$ c.p.d. $\Rightarrow$ $\int_{E} g \leq \int_{E} f$; y $E_{1} \subseteq E_{2}$ $\Rightarrow$ $\int_{E_{1}} f \leq \int_{E_{2}} f$. \hfill\textit{3.3.3 y 4.4.2.}
\item \textbf{Linealidad:} $\int_{E} (cf + g) = c \int_{E} f + \int_{E} g$. \hfill\textit{Teorema 4.5.1.}
\item \textbf{Aditividad del dominio:} $E = \bigcup_{k} E_{k}$ disjuntos $\Rightarrow$ $\int_{E} f = \sum_{k} \int_{E_{k}} f$. \hfill\textit{Proposición 4.2.1.}
\item \textbf{Markov:} $m(\{f \geq \alpha\}) \leq \frac{1}{\alpha} \int_{E} f \, dx$ para $\alpha > 0$. \hfill\textit{Proposición 4.4.3.}
\item \textbf{Anulación:} $\int_{E} f = 0 \iff f = 0$ c.p.d.; $m(E) = 0 \Rightarrow \int_{E} f = 0$. \hfill\textit{Corolario 4.4.4 y Lema 4.4.1.}
\item \textbf{Finitud:} $\int_{E} f < \infty \Rightarrow f < \infty$ c.p.d. \hfill\textit{Proposición 4.1.3.}
\end{itemize}
\end{equivbox}

\anexogroup{La integral general y $L(E)$}

\begin{equivbox}{brown!85!black}{La integral de funciones con signo}
Sea $f : E \to \overline{\R}$ medible, $f = f^{+} - f^{-}$.
\begin{itemize}[leftmargin=2em,itemsep=2pt]
\item $\int_{E} f$ \emph{existe} si $\int f^{+} < \infty$ o $\int f^{-} < \infty$; $f \in L(E)$ si ambas son finitas. \hfill\textit{5.1.1 y 5.1.2.}
\item $f \in L(E) \iff \int_{E} |f| \, dx < \infty$, y en tal caso $\big| \int_{E} f \big| \leq \int_{E} |f|$ y $f$ es finita c.p.d. \hfill\textit{Nota 5.1.3.}
\item \textbf{Monotonía:} $f \leq g$ c.p.d. y ambas integrales existen $\Rightarrow$ $\int_{E} f \leq \int_{E} g$. \hfill\textit{Teorema 5.2.1.}
\item \textbf{Linealidad en $L(E)$:} $\int cf = c \int f$; $\int (f+g) = \int f + \int g$. \hfill\textit{Teorema 5.3.1.}
\item \textbf{Resta con minorante:} $\phi \leq f$ c.p.d., $\phi \in L(E)$ $\Rightarrow$ $\int (f - \phi) = \int f - \int \phi$. \hfill\textit{Proposición 5.3.3.}
\item \textbf{Cuidado:} $\int (f - g) = \int f - \int g$ puede fallar si las integrales no son finitas. \hfill\textit{Nota 5.3.2.}
\item \textbf{Productos:} $f \in L(E)$, $|g| \leq M$ medible $\Rightarrow fg \in L(E)$. \hfill\textit{Nota 5.3.4.}
\end{itemize}
\end{equivbox}

\anexogroup{Riemann y Lebesgue}

\begin{equivbox}{red!75!black}{Criterio de Lebesgue para la integrabilidad de Riemann}
Sea $f : [a,b] \to \R$ acotada. Las siguientes son equivalentes (Teorema 6.2.1):
\begin{enumerate}[label=(\arabic*),leftmargin=2em]
\item \textbf{(Riemann)}\; $f$ es Riemann integrable.
\item \textbf{(continuidad c.p.d.)}\; El conjunto de discontinuidades de $f$ tiene medida cero.
\end{enumerate}
En tal caso, la integral de Riemann de $f$ coincide con su integral de Lebesgue $\int_{[a,b]} f \, dx$.
\par\medskip
\textbf{Mecanismo:} escalonadas $\ell_{k} \leq f \leq u_{k}$ sobre particiones refinadas con malla $\to 0$; TCD da $\int \ell_{k} \to \int \ell$, $\int u_{k} \to \int u$; y $f$ Riemann integrable $\iff u = \ell = f$ c.p.d. (Proposición 6.1.3 y Nota 6.1.4).
\par\medskip
\textbf{Contraejemplo canónico:} la función de Dirichlet $\ind_{\Q \cap [0,1]}$ es discontinua en todo punto: no es Riemann integrable, pero es Lebesgue integrable con $\int_{[0,1]} \ind_{\Q} \, dx = 0$.
\end{equivbox}

\newpage
% =============================================================================
\section{Anexo: contraejemplos canónicos}
% =============================================================================

Esta tabla recopila los contraejemplos más útiles para el examen --- cada uno responde a una pregunta del estilo ``¿es necesariamente \textit{X}?'' con \textit{no, esta sucesión/función cumple A pero no B}.

\renewcommand{\arraystretch}{1.35}
\begin{longtable}{p{0.30\textwidth} p{0.32\textwidth} p{0.30\textwidth}}
\hline
\textbf{Propiedad afirmativa} & \textbf{Contraejemplo a la implicación} & \textbf{Comentario / referencia} \\
\hline
\endhead

Convergencia c.p.d. $\Rightarrow$ casi uniforme (Egorov) sin $m(E) < \infty$ & $f_{k} = \ind_{B(0,k)} \to 1$ puntualmente en $\R^{d}$; no hay convergencia uniforme en ningún cerrado con complemento de medida finita. & La medida finita es esencial en Egorov (Ejemplo 1.1.2). \\
\hline
Convergencia c.p.d. $\Rightarrow$ en medida & $f_{k} = \ind_{[k,k+1]} \to 0$ puntualmente en $\R$, pero $m\{|f_{k}| > \tfrac{1}{2}\} = 1$ para todo $k$. & Requiere $m(E) < \infty$ (Teorema 2.1.3). \\
\hline
Convergencia en medida $\Rightarrow$ c.p.d. & Máquina de escribir: indicadoras diádicas $I_{n,i}$ en $[0,1]$; $\xrightarrow{m} 0$ pero no converge en ningún punto. & Solo se rescata una subsucesión c.p.d. (Ejemplo 2.1.4, Teorema 2.1.5). \\
\hline
$\lim_{k} \int f_{k} = \int \lim_{k} f_{k}$ sin hipótesis extra & $f_{k} = k \ind_{[0,1/k]}$ en $[0,1]$: $\int f_{k} = 1$ pero $f_{k} \to 0$ c.p.d. & Falta dominación: $\sup_{k} f_{k}(x) \approx \tfrac{1}{x} \notin L([0,1])$ (Ejemplo 5.4.3). \\
\hline
Desigualdad de Fatou con igualdad & $f_{k} = \ind_{[k,k+1]}$: $\int \liminf f_{k} = 0 < 1 = \liminf \int f_{k}$. & Fatou solo da ``$\leq$'' (Teorema 5.4.4). \\
\hline
TCM decreciente sin mayorante integrable & $f_{k} = \ind_{[k,\infty)} \downarrow 0$ puntualmente, pero $\int_{\R} f_{k} = \infty \not\to 0$. & En 5.4.1(ii) la mayorante $\phi \in L(E)$ es esencial. \\
\hline
Convergencia uniforme $\Rightarrow$ $\lim \int = \int \lim$ sin $m(E) < \infty$ & $f_{k} = \tfrac{1}{k} \to 0$ uniformemente en $\R$, pero $\int_{\R} f_{k} = \infty$. & La medida finita es esencial en 5.4.7 (Nota 5.4.8). \\
\hline
$\int (f - g) = \int f - \int g$ sin integrabilidad & $f = \ind_{[n,\infty)}$, $g = \ind_{[n+1,\infty)}$: $\int (f-g) = 1$, pero $\int f - \int g = \infty - \infty$. & La linealidad requiere $f, g \in L(E)$ (Nota 5.3.2). \\
\hline
Acotada $\Rightarrow$ Riemann integrable & Dirichlet $\ind_{\Q \cap [0,1]}$: acotada, Lebesgue integrable con integral $0$, pero discontinua en todo punto. & Criterio de Lebesgue (Teorema 6.2.1). \\
\hline
Integral finita $\Rightarrow$ función finita en todo punto & $f = \infty \cdot \ind_{Z}$ con $m(Z) = 0$: $\int f = 0$ pero $f = \infty$ en $Z$. & Solo se obtiene finitud c.p.d. (Proposición 4.1.3). \\
\hline
$\int_{E} f = 0 \Rightarrow f = 0$ en todo punto & $f = \ind_{Z}$ con $m(Z) = 0$, $Z \neq \emptyset$. & Solo $f = 0$ c.p.d. (Corolario 4.4.4). \\
\hline
\end{longtable}
\renewcommand{\arraystretch}{1.0}
\newpage
% =============================================================================
\section{Anexo: los teoremas de convergencia en un vistazo}
% =============================================================================

Mapa rápido para decidir qué teorema permite intercambiar $\lim$ e $\int$. Cada bloque indica hipótesis, conclusión y qué falla si se quita una hipótesis (ver \emph{Anexo: contraejemplos canónicos}).

\medskip

\renewcommand{\arraystretch}{1.35}
\begin{longtable}{p{0.22\textwidth} p{0.36\textwidth} p{0.34\textwidth}}
\hline
\textbf{Teorema} & \textbf{Hipótesis} & \textbf{Conclusión} \\
\hline
\endhead

TCM no negativo (4.2.2) & $0 \leq f_{k} \leq f_{k+1}$ en $E$; $f_{k} \to f$ c.p.d. & $\int_{E} f_{k} \, dx \to \int_{E} f \, dx$ (posiblemente $= \infty$). \\
\hline
TCM con minorante / mayorante (5.4.1) & $\phi \in L(E)$; $\phi \leq f_{k} \leq f_{k+1}$ c.p.d. (o bien $f_{k+1} \leq f_{k} \leq \phi$); $f_{k} \to f$ c.p.d. & $\int_{E} f_{k} \, dx \to \int_{E} f \, dx$. \\
\hline
Series no negativas (5.4.2) & $f_{k} \geq 0$ medibles. & $\int_{E} \sum_{k} f_{k} \, dx = \sum_{k} \int_{E} f_{k} \, dx$. \\
\hline
Lema de Fatou (5.4.4) & $f_{k} \geq 0$ medibles. & $\int_{E} \liminf f_{k} \, dx \leq \liminf \int_{E} f_{k} \, dx$ (solo desigualdad). \\
\hline
TCD (5.4.6) & $0 \leq f_{k} \leq \phi$ c.p.d.; $\int_{E} \phi \, dx < \infty$; $f_{k} \to f$ c.p.d. & $\int_{E} f_{k} \, dx \to \int_{E} f \, dx$. \\
\hline
Convergencia uniforme (5.4.7) & $f_{k} \in L(E)$; $f_{k} \to f$ \emph{uniformemente}; $m(E) < \infty$. & $f \in L(E)$ y $\int_{E} f_{k} \, dx \to \int_{E} f \, dx$ (prueba: ejercicio). \\
\hline
\end{longtable}
\renewcommand{\arraystretch}{1.0}

\begin{equivbox}{olive!80!black}{Implicaciones y advertencias alrededor de la convergencia}
\textbf{Implicaciones siempre válidas.}
\begin{itemize}[leftmargin=2em,itemsep=2pt]
\item Monotonía creciente $+$ no negatividad $\Rightarrow$ vale TCM sin más hipótesis. \hfill\textit{Teorema 4.2.2.}
\item Dominación integrable $\Rightarrow$ vale TCD. \hfill\textit{Teorema 5.4.6.}
\item Solo no negatividad $\Rightarrow$ vale Fatou (desigualdad). \hfill\textit{Teorema 5.4.4.}
\item $\int_{E} f_{k} \, dx \leq M$ para todo $k$ $\Rightarrow$ $\int_{E} \liminf f_{k} \, dx \leq M$. \hfill\textit{Nota 5.4.5.}
\end{itemize}
\textbf{Implicaciones que \emph{no} valen en general.}
\begin{itemize}[leftmargin=2em,itemsep=2pt]
\item Convergencia puntual $\not\Rightarrow$ $\lim \int = \int \lim$ (e.g., $k \ind_{[0,1/k]}$; $\ind_{[k,k+1]}$).
\item Fatou con igualdad $\not\Rightarrow$ (e.g., $\ind_{[k,k+1]}$: $0 < 1$).
\item Convergencia uniforme $\not\Rightarrow$ $\lim \int = \int \lim$ si $m(E) = \infty$ (e.g., $f_{k} = \tfrac{1}{k}$ en $\R$).
\item Monotonía decreciente sin mayorante integrable $\not\Rightarrow$ TCM (e.g., $\ind_{[k,\infty)}$).
\end{itemize}
\end{equivbox}

\newpage
% =============================================================================
\section{Anexo: estrategias de demostración}
% =============================================================================

Patrones meta de demostración que recurren a lo largo de este parcial. Cuando un problema de examen pida ``muestre que \dots'', revise estos patrones antes de empezar a escribir.

\newtcolorbox{strategybox}[1]{%
  enhanced,
  breakable,
  colback=gray!4,
  colframe=blue!50!black,
  coltitle=white,
  fonttitle=\bfseries,
  fontupper=\small,
  title={#1},
  arc=3pt,
  boxrule=0.6pt,
  left=8pt, right=8pt, top=4pt, bottom=4pt,
  before skip=8pt, after skip=8pt,
}

\begin{strategybox}{1.\ Para probar que algo vale c.p.d.: hacer pequeño el conjunto malo}
\textbf{Patrón:} \textit{exprese el conjunto malo como unión numerable de conjuntos que puede controlar.}
\begin{itemize}[leftmargin=1.5em]
\item $\{f > 0\} = \bigcup_{k} \{f \geq \tfrac{1}{k}\}$ y controle cada pieza con Markov.
\item Conjuntos $H_{m} = \bigcup_{j \geq m} A_{j}$ decrecientes con $m(H_{m}) \leq \sum_{j \geq m} m(A_{j}) \to 0$ cuando $\sum m(A_{j}) < \infty$ (estilo Borel-Cantelli); el conjunto malo es $\bigcap_{m} H_{m}$.
\item Presupueste con series geométricas: pida $m(\cdot) < \varepsilon / 2^{k}$ en el paso $k$.
\end{itemize}
\textbf{Apariciones:} subsucesión c.p.d. de convergencia en medida (§2.1), $\int f = 0 \Rightarrow f = 0$ c.p.d. (§4.4), Egorov (§1.1).
\end{strategybox}

\begin{strategybox}{2.\ Para intercambiar límite e integral: lista de chequeo}
\textbf{Patrón:} \textit{identifique cuál teorema de convergencia encaja con los datos.}
\begin{itemize}[leftmargin=1.5em]
\item ¿La sucesión es monótona y no negativa (o con minorante/mayorante en $L(E)$)? $\Rightarrow$ TCM.
\item ¿Hay una dominante $\phi$ con $\int \phi < \infty$? $\Rightarrow$ TCD.
\item ¿Solo necesita una desigualdad? $\Rightarrow$ Fatou (recuerde: $\liminf$ adentro $\leq$ $\liminf$ afuera).
\item ¿Convergencia uniforme y $m(E) < \infty$? $\Rightarrow$ teorema de convergencia uniforme.
\item Si nada aplica, sospeche del resultado: busque contraejemplo tipo $k \ind_{[0,1/k]}$ o $\ind_{[k,k+1]}$.
\end{itemize}
\textbf{Apariciones:} sumas de Darboux (§6.1), series término a término (§5.4), TCD vía Fatou (§5.4).
\end{strategybox}

\begin{strategybox}{3.\ Para probar propiedades de la integral: simple $\to$ no negativa $\to$ general}
\textbf{Patrón:} \textit{pruebe primero para funciones simples, suba por TCM, y separe $f = f^{+} - f^{-}$.}
\begin{itemize}[leftmargin=1.5em]
\item Para simples: use la partición común $\{A_{i} \cap B_{j}\}$ y la aditividad finita de $m$.
\item Para no negativas: tome simples $\phi_{k} \uparrow f$ (aproximación diádica) y aplique TCM.
\item Para generales: aplique lo anterior a $f^{+}$ y $f^{-}$ y reste, cuidando que no aparezca $\infty - \infty$.
\end{itemize}
\textbf{Apariciones:} linealidad (§4.5 y §5.3), monotonía general (§5.2), homogeneidad (§4.5).
\end{strategybox}

\begin{strategybox}{4.\ Para construir cerrados donde pasa algo bueno: regularidad $+$ presupuesto geométrico}
\textbf{Patrón:} \textit{regularidad interna da $F_{k}$ cerrado dentro de cada pieza con $m(\cdot \setminus F_{k}) < \varepsilon/2^{k}$; luego combine.}
\begin{itemize}[leftmargin=1.5em]
\item Intersección numerable $\bigcap F_{k}$: sigue siendo cerrada y pierde a lo sumo $\sum \varepsilon/2^{k} = \varepsilon$ de medida (Egorov, Lusin caso finito).
\item Unión numerable $\bigcup F_{k}$: no es cerrada en general, pero sí lo es si los $F_{k}$ viven en anillos $\{k - 1 \leq |x| < k\}$ (localmente finita): toda sucesión convergente cae en finitos $F_{k}$ (Lusin caso infinito).
\item Para continuidad relativa a una unión finita de cerrados disjuntos: una sucesión convergente eventualmente permanece en el cerrado que contiene al límite (§1.2).
\end{itemize}
\textbf{Apariciones:} lema técnico de Egorov (§1.1), propiedad $C$ de simples y Lusin (§1.2).
\end{strategybox}

\begin{strategybox}{5.\ Para pasar de medida finita a infinita (y de acotado a no acotado)}
\textbf{Patrón:} \textit{trocee el dominio y use continuidad de la medida.}
\begin{itemize}[leftmargin=1.5em]
\item Anillos: $E_{k} = E \cap \{k-1 \leq |x| < k\}$, disjuntos, de medida finita, con $E = \bigcup_{k} E_{k}$.
\item Truncamiento creciente: $A_{k} = A \cap B(0,k) \uparrow A$ y $m(A) = \lim m(A_{k})$ (continuidad desde abajo).
\item Reparta el error: pida $\varepsilon/2^{k}$ en la pieza $k$-ésima.
\end{itemize}
\textbf{Apariciones:} Lusin (§1.2), medida de $A \times [0,a]$ (§3.1), gráfico de medida cero (§3.2, ejercicio).
\end{strategybox}

\begin{strategybox}{6.\ Para trabajar con $R(f,E)$: geometría de la integral}
\textbf{Patrón:} \textit{traduzca afirmaciones sobre integrales a afirmaciones sobre medidas de regiones.}
\begin{itemize}[leftmargin=1.5em]
\item Monotonía de la integral $=$ monotonía de la medida: $g \leq f \Rightarrow R(g,E) \subseteq R(f,E)$.
\item Aditividad del dominio $=$ aditividad de la medida: dominios disjuntos dan regiones disjuntas.
\item TCM $=$ continuidad desde abajo: $R(f_{k},E) \uparrow$ y $R(f,E) = \Gamma(f,E) \cup \bigcup_{k} R(f_{k},E)$ con $m(\Gamma) = 0$.
\item Restricción: $\int_{E_{1}} f = \int_{E} \ind_{E_{1}} f$ porque las regiones difieren en un conjunto nulo.
\end{itemize}
\textbf{Apariciones:} definición de la integral y sus propiedades (§3.3, §4.2).
\end{strategybox}

\begin{strategybox}{7.\ Para usar Egorov / Lusin en pruebas de convergencia}
\textbf{Patrón:} \textit{cambie ``c.p.d.'' por ``uniforme en un cerrado grande'' y pague $\eta$ de medida.}
\begin{itemize}[leftmargin=1.5em]
\item Egorov: si $m(E) < \infty$ y $f_{n} \to f$ c.p.d., hay $F$ cerrado con $m(E \setminus F) < \eta$ y convergencia uniforme en $F$.
\item Con la uniformidad, los conjuntos $\{|f_{n} - f| > \varepsilon\}$ acaban dentro de $E \setminus F$: así se prueba la convergencia en medida (§2.1).
\item Lusin: cambie ``$f$ medible'' por ``$f$ continua en un cerrado grande'' cuando necesite continuidad.
\end{itemize}
\textbf{Apariciones:} c.p.d. $\Rightarrow$ en medida (§2.1), Lusin desde Egorov (§1.2).
\end{strategybox}

\newpage
% =============================================================================
\section{Anexo: resumen de fórmulas}
% =============================================================================

\begin{equivbox}{orange}{Resumen rápido --- todas las fórmulas en un vistazo}
Sean $f, g : E \to [0,\infty]$ (o $\overline{\R}$ donde tenga sentido) medibles, $c \in \R$, $\alpha > 0$, $\phi$ simple.
\begin{itemize}[leftmargin=2em]
\item \textbf{Definición:} $\int_{E} f \, dx = m(R(f,E))$ para $f \geq 0$; $\int_{E} f \, dx = \int_{E} f^{+} \, dx - \int_{E} f^{-} \, dx$ en general.
\item \textbf{Simples:} $\int_{E} \sum_{i} a_{i} \ind_{A_{i}} \, dx = \sum_{i} a_{i}\, m(A_{i})$ ($A_{i}$ disjuntos).
\item \textbf{Producto:} $m(A \times [0,a]) = a\, m(A)$; $m_{e}(\Gamma(f,E)) = 0$.
\item \textbf{Supremos:} $\int_{E} f \, dx = \sup \sum_{i} (\inf_{E_{i}} f)\, m(E_{i}) = \sup \{ \int_{E} \phi \, dx : 0 \leq \phi \leq f \}$.
\item \textbf{Restricción:} $\int_{E_{1}} f \, dx = \int_{E} \ind_{E_{1}} f \, dx$; dominio nulo $\Rightarrow$ integral nula.
\item \textbf{Aditividad:} $E = \bigcup_{k} E_{k}$ disjuntos $\Rightarrow$ $\int_{E} f = \sum_{k} \int_{E_{k}} f$.
\item \textbf{Linealidad:} $\int (cf + g) = c \int f + \int g$ (para $f, g \geq 0$ con $c \geq 0$, o $f, g \in L(E)$).
\item \textbf{Markov:} $m(\{f \geq \alpha\}) \leq \frac{1}{\alpha} \int_{E} f \, dx$.
\item \textbf{Anulación:} $\int_{E} |f| \, dx = 0 \iff f = 0$ c.p.d.
\item \textbf{Triangular:} $|\int_{E} f \, dx| \leq \int_{E} |f| \, dx$; $f \in L(E) \iff \int |f| < \infty$.
\item \textbf{TCM:} $0 \leq f_{k} \uparrow f$ c.p.d. $\Rightarrow$ $\int f_{k} \to \int f$; con minorante $\phi \in L(E)$ también.
\item \textbf{Fatou:} $\int \liminf f_{k} \leq \liminf \int f_{k}$ ($f_{k} \geq 0$).
\item \textbf{TCD:} $0 \leq f_{k} \leq \phi$, $\int \phi < \infty$, $f_{k} \to f$ c.p.d. $\Rightarrow$ $\int f_{k} \to \int f$.
\item \textbf{Series:} $f_{k} \geq 0 \Rightarrow \int \sum_{k} f_{k} = \sum_{k} \int f_{k}$.
\item \textbf{En medida:} $f_{n} \xrightarrow{m} f \iff m\{|f_{n} - f| > \varepsilon\} \to 0\ \forall \varepsilon > 0$; c.p.d. $+$ $m(E) < \infty$ $\Rightarrow$ en medida; en medida $\Rightarrow$ subsucesión c.p.d.
\item \textbf{Egorov:} $m(E) < \infty$ $\Rightarrow$ convergencia c.p.d. es uniforme fuera de un cerrado de medida $< \varepsilon$.
\item \textbf{Lusin:} $f$ medible $\iff$ $f$ continua relativa a un cerrado $F$ con $m(E \setminus F) < \varepsilon$, $\forall \varepsilon$.
\item \textbf{Riemann:} $f$ acotada en $[a,b]$ es Riemann integrable $\iff$ continua c.p.d.; en tal caso ambas integrales coinciden.
\end{itemize}
\end{equivbox}

\end{document}
